\documentclass[../main.tex]{subfiles}
\begin{document}

\chapter{GPU Architecture Optimizations - II}
\lessoninfo{
  video={M7, lessons 1--3 (GPU Virtual Memory; GPU Virtual Memory - UVA; GPU Warp Scheduling)},
  sources={Power et al.\ \cite{power2014x86} (GPU MMU); Ausavarungnirun et al.\ \cite{ausavarungnirun2017mosaic} (multiple page sizes); Li et al.\ \cite{li2021tlb} (multi-GPU TLBs); Ganguly et al.\ \cite{ganguly2019interplay} (prefetching and eviction in unified memory); Kim et al.\ \cite{kim2020batch} (batched fault handling); Rogers et al.\ \cite{rogers2012ccws} (cache-conscious wavefront scheduling); Narasiman et al.\ \cite{narasiman2011twolevel} and Gebhart et al.\ \cite{gebhart2011energy} (two-level scheduling); Jog et al.\ \cite{jog2013orchestrated} (prefetch-aware scheduling), \cite{jog2013owl} (OWL)},
  textbook={A18 \cite{aamodt2018gpu} ch.~3, ch.~4; K12 \cite{kim2012gpgpu} §1.4.1 GPU Pipeline; CPG \cite{nvidia2024cudaguide} §4.2 Hardware Multithreading, §19 Unified Memory Programming; BPG \cite{nvidia2026cudabestpractices} §10.1.4 Unified Virtual Addressing},
  keywords={virtual memory, address translation, TLB, multi-level page table, page table walker, large page, segment mapping, copy-then-execute, shared virtual memory, on-demand paging, far fault, tree-based neighborhood prefetching, oversubscription, warp scheduler, round-robin, greedy-then-oldest, two-level scheduling, working set, cache-conscious wavefront scheduling},
}

Lesson~6 changed the inside of the SM: how a warp survives a divergent
branch and how its registers are stored.  This lesson continues with two
further parts of the SM's machinery.  The first two sections concern
memory addresses.  A GPU that shares data structures with the CPU needs
virtual memory, and translating the addresses of a whole warp at once
strains the translation hardware in a way a CPU core never does; when
CPU and GPU share one address space, pages also have to be allocated
and moved on demand, across PCIe, while the kernel runs.  The third
section returns to the scheduler stage of the pipeline of Lesson~4 and
asks which warp should issue when several are ready, a choice that
decides how well memory latency is hidden and how much of the L1 cache
each warp can keep.  Like Lesson~6, the lesson is organised around
research papers; each section names the ones it draws on.

\section{GPU virtual memory}
\label{sec:l07-vm}

Virtual memory has two sides: the translation of every address the
program issues, and the allocation of the physical pages behind
them.\source{M7, GPU virtual memory; A18 ch.~4; \cite{power2014x86};
CPG §19.2.2.1 Memory Paging and Page Sizes.}  On a CPU the operating
system allocates pages.  On a GPU the allocation was originally done by
the CPU together with the GPU driver, through \texttt{cudaMalloc},
\texttt{cudaMallocManaged} and \texttt{cudaMemcpy}; \cref{sec:l07-uva}
returns to it.  This section reviews address translation, shows what
is different when a warp rather than a thread issues the addresses, and
presents three ways to make translation cheaper.

\subsection{Address translation}

\begin{definition}[Virtual memory]
  \label{def:l07-vm}
  With \term{virtual memory}[仮想メモリ] a program issues
  \term{virtual address}[仮想アドレス]es in a private address space,
  and the hardware maps each one, at every access, to a
  \term{physical address}[物理アドレス] in the memory.  The mapping is
  called \term{address translation}[アドレス変換].  It is kept at the
  granularity of a \term{page}[ページ], a fixed-size, aligned block of
  the address space, typically \SI{4}{\kibi\byte} or
  \SI{2}{\mebi\byte}.
\end{definition}

A virtual address splits into two fields (\cref{fig:l07-translation}).
The upper bits are the
\term{virtual page number (VPN)}[仮想ページ番号], the lower bits the
\term{page offset}[ページオフセット], the position inside the page.
Translation replaces the VPN by a
\term{physical frame number (PFN)}[物理フレーム番号], also called the
physical page number (PPN), and keeps the offset unchanged; a page of
$2^k$ bytes has a $k$-bit offset.\index{PFN|see{physical frame number (PFN)}}\index{VPN|see{virtual page number (VPN)}}\index{PPN|see{physical frame number (PFN)}}

\begin{marginfigure}
  % Address translation (Lesson 7): a virtual address splits into a
% virtual page number and a page offset; translation replaces the VPN
% by a physical frame number and keeps the offset unchanged.  Field
% widths are those of 4 KiB pages with a 48-bit virtual address.
% Redraws the outline figure m7l1_translation.
% Usage: % Address translation (Lesson 7): a virtual address splits into a
% virtual page number and a page offset; translation replaces the VPN
% by a physical frame number and keeps the offset unchanged.  Field
% widths are those of 4 KiB pages with a 48-bit virtual address.
% Redraws the outline figure m7l1_translation.
% Usage: % Address translation (Lesson 7): a virtual address splits into a
% virtual page number and a page offset; translation replaces the VPN
% by a physical frame number and keeps the offset unchanged.  Field
% widths are those of 4 KiB pages with a 48-bit virtual address.
% Redraws the outline figure m7l1_translation.
% Usage: \input{figures/l07-translation}  (in a marginfigure)
% Width: margin (50 mm); the drawing is about 46 mm wide.
\begin{tikzpicture}[x=1mm, y=1mm, font=\scriptsize\sffamily,
    >={Stealth[length=1.6mm]}]
  % virtual address
  \node[anchor=east] at (-0.5,23) {VA};
  \draw[fill=ln-accent!20] (0,20) rectangle (28,26) node[midway] {VPN};
  \draw[fill=ln-box] (28,20) rectangle (44,26) node[midway] {\iflangja{オフセット}{offset}};
  \node[anchor=south, text=ln-muted] at (0.8,26) {47};
  \node[anchor=south, text=ln-muted] at (28,26) {12\,11};
  \node[anchor=south, text=ln-muted] at (43.2,26) {0};
  % translation
  \draw[->] (14,20) -- (14,13) node[midway, right, align=left, text=ln-muted]
    {TLB /\\\iflangja{ページテーブル}{page table}};
  \draw[->] (36,20) -- (36,13) node[midway, right, text=ln-muted]
    {\iflangja{不変}{copied}};
  % physical address
  \node[anchor=east] at (-0.5,10) {PA};
  \draw[fill=ln-tip!25] (0,7) rectangle (28,13) node[midway] {PFN};
  \draw[fill=ln-box] (28,7) rectangle (44,13) node[midway] {\iflangja{オフセット}{offset}};
\end{tikzpicture}
  (in a marginfigure)
% Width: margin (50 mm); the drawing is about 46 mm wide.
\begin{tikzpicture}[x=1mm, y=1mm, font=\scriptsize\sffamily,
    >={Stealth[length=1.6mm]}]
  % virtual address
  \node[anchor=east] at (-0.5,23) {VA};
  \draw[fill=ln-accent!20] (0,20) rectangle (28,26) node[midway] {VPN};
  \draw[fill=ln-box] (28,20) rectangle (44,26) node[midway] {\iflangja{オフセット}{offset}};
  \node[anchor=south, text=ln-muted] at (0.8,26) {47};
  \node[anchor=south, text=ln-muted] at (28,26) {12\,11};
  \node[anchor=south, text=ln-muted] at (43.2,26) {0};
  % translation
  \draw[->] (14,20) -- (14,13) node[midway, right, align=left, text=ln-muted]
    {TLB /\\\iflangja{ページテーブル}{page table}};
  \draw[->] (36,20) -- (36,13) node[midway, right, text=ln-muted]
    {\iflangja{不変}{copied}};
  % physical address
  \node[anchor=east] at (-0.5,10) {PA};
  \draw[fill=ln-tip!25] (0,7) rectangle (28,13) node[midway] {PFN};
  \draw[fill=ln-box] (28,7) rectangle (44,13) node[midway] {\iflangja{オフセット}{offset}};
\end{tikzpicture}
  (in a marginfigure)
% Width: margin (50 mm); the drawing is about 46 mm wide.
\begin{tikzpicture}[x=1mm, y=1mm, font=\scriptsize\sffamily,
    >={Stealth[length=1.6mm]}]
  % virtual address
  \node[anchor=east] at (-0.5,23) {VA};
  \draw[fill=ln-accent!20] (0,20) rectangle (28,26) node[midway] {VPN};
  \draw[fill=ln-box] (28,20) rectangle (44,26) node[midway] {\iflangja{オフセット}{offset}};
  \node[anchor=south, text=ln-muted] at (0.8,26) {47};
  \node[anchor=south, text=ln-muted] at (28,26) {12\,11};
  \node[anchor=south, text=ln-muted] at (43.2,26) {0};
  % translation
  \draw[->] (14,20) -- (14,13) node[midway, right, align=left, text=ln-muted]
    {TLB /\\\iflangja{ページテーブル}{page table}};
  \draw[->] (36,20) -- (36,13) node[midway, right, text=ln-muted]
    {\iflangja{不変}{copied}};
  % physical address
  \node[anchor=east] at (-0.5,10) {PA};
  \draw[fill=ln-tip!25] (0,7) rectangle (28,13) node[midway] {PFN};
  \draw[fill=ln-box] (28,7) rectangle (44,13) node[midway] {\iflangja{オフセット}{offset}};
\end{tikzpicture}

  \caption{Translation keeps the page offset and replaces the VPN by a
    PFN (bit positions for \SI{4}{\kibi\byte} pages and a 48-bit
    virtual address).}
  \label{fig:l07-translation}
\end{marginfigure}

\begin{definition}[Page table, TLB and page fault]
  \label{def:l07-tlb}
  The \term{page table}[ページテーブル] is the data structure, kept in
  memory, that holds the mapping; each
  \term{page table entry (PTE)}[ページテーブルエントリ]\index{PTE|see{page table entry (PTE)}} maps one VPN to
  a PFN and carries metadata such as valid and permission bits.  A
  \term{TLB (translation lookaside buffer)}[TLB]\index{translation lookaside buffer|see{TLB (translation lookaside buffer)}} is a small cache of
  recently used translations.  A \term{TLB miss}[TLBミス] occurs when
  the TLB does not hold the translation, so the page table must be read;
  a \term{page fault}[ページフォールト] occurs when the page table holds
  no valid mapping because no physical page has been allocated yet, so
  that the operating system must allocate one.
\end{definition}

\caution{A TLB miss and a page fault are different events.  A TLB miss
is resolved by hardware (or a short handler) reading the page table; a
page fault needs the operating system and, for a GPU, the host
(\cref{sec:l07-uva}).}%
The TLB exists because the PTE itself is in memory: without it, every
load would first need one or more memory accesses to find its
translation.

\subsection{Multi-level page tables}

A single table indexed by the VPN would need $2^{v}$ entries for a
$v$-bit VPN, whether or not the program uses the pages.

\begin{definition}[Multi-level page table and page table walk]
  \label{def:l07-multilevel}
  A \term{multi-level page table}[多段ページテーブル] splits the VPN
  into several indices, one per level.  The index of the top level
  selects an entry of the top-level table that holds the base address
  of a next-level table, whose entry is selected by the next index, and
  so on; the entry of the last level is the PTE.  Following this chain
  is a \term{page table walk}[ページテーブルウォーク].  It is performed
  either by a \term{page table walker (PTW)}[ページテーブルウォーカ], a
  hardware state machine, or by a software miss handler.  Tables that
  no used address reaches are never allocated.\index{PTW|see{page table walker (PTW)}}
\end{definition}

\Cref{fig:l07-page-walk} shows the four-level form used by x86-64,
ARM and GPUs.\margin{x86-64 with four levels: 48-bit virtual
addresses, four 9-bit indices and a 12-bit offset; each table has 512
entries of \SI{8}{\byte}, exactly one \SI{4}{\kibi\byte} page; the
root register is CR3.}  The
walk starts at a root register, uses the level-4
index to read the level-3 table base, and so on down to the level-1
entry, the PTE; PFN and offset together form the physical address.
The price of the saved space is time: one translation that misses in
the TLB costs at least four dependent memory accesses.

\begin{figure}[tb]
  \centering
  % Four-level page table walk (Lesson 7), x86-64 layout: a 48-bit
% virtual address holds four 9-bit indices and a 12-bit page offset.
% CR3 points to the level-4 table; each index selects an entry that
% holds the base of the next table; the level-1 entry is the PTE whose
% PFN, joined with the offset, forms the physical address.
% Redraws the outline figure m7l1_fourleveltlb (bit widths re-derived
% for x86-64, slide erratum of outline Appendix C).
% Usage: % Four-level page table walk (Lesson 7), x86-64 layout: a 48-bit
% virtual address holds four 9-bit indices and a 12-bit page offset.
% CR3 points to the level-4 table; each index selects an entry that
% holds the base of the next table; the level-1 entry is the PTE whose
% PFN, joined with the offset, forms the physical address.
% Redraws the outline figure m7l1_fourleveltlb (bit widths re-derived
% for x86-64, slide erratum of outline Appendix C).
% Usage: % Four-level page table walk (Lesson 7), x86-64 layout: a 48-bit
% virtual address holds four 9-bit indices and a 12-bit page offset.
% CR3 points to the level-4 table; each index selects an entry that
% holds the base of the next table; the level-1 entry is the PTE whose
% PFN, joined with the offset, forms the physical address.
% Redraws the outline figure m7l1_fourleveltlb (bit widths re-derived
% for x86-64, slide erratum of outline Appendix C).
% Usage: \input{figures/l07-page-walk}
% Width: body (120 mm); the drawing is about 116 mm wide.
\begin{tikzpicture}[x=0.95mm, y=1mm, font=\scriptsize\sffamily,
    >={Stealth[length=1.6mm]}]
  % virtual address fields
  \node[anchor=east] at (-1,47) {VA};
  \foreach \i/\lab/\bits in {0/L4/47--39, 1/L3/38--30, 2/L2/29--21, 3/L1/20--12} {
    \draw[fill=ln-accent!20] (\i*22,44) rectangle ++(22,6)
      node[midway] {\lab\ \iflangja{インデックス}{index}};
    \node[anchor=south, text=ln-muted] at (\i*22+11,50) {\bits};
  }
  \draw[fill=ln-box] (88,44) rectangle (112,50)
    node[midway] {\iflangja{オフセット}{offset}};
  \node[anchor=south, text=ln-muted] at (100,50) {11--0};
  % tables: three rows each, the middle one selected
  \foreach \i/\lab/\entry in {0/L4/{L3 base}, 1/L3/{L2 base}, 2/L2/{L1 base}, 3/L1/{PTE: PFN}} {
    \draw[fill=white] (\i*22+3,32) rectangle ++(16,-4) node[midway] {$\cdots$};
    \draw[fill=ln-tip!25] (\i*22+3,28) rectangle ++(16,-5) node[midway] {\entry};
    \draw[fill=white] (\i*22+3,23) rectangle ++(16,-4) node[midway] {$\cdots$};
    \node[text=ln-muted] at (\i*22+11,16.5) {\lab\ \iflangja{テーブル}{table}};
    \draw[->] (\i*22+11,44) -- (\i*22+11,32);
  }
  % next-table pointers
  \foreach \i in {0,1,2} {
    \draw[->, ln-accent, thick] (\i*22+19,25.5) -- ++(2,0) |- (\i*22+25,31);
  }
  % CR3
  \draw[->, ln-accent, thick] (-2,30) node[left, text=black] {CR3} -- (3,30);
  % physical address
  \node[anchor=east] at (61,6) {PA};
  \draw[fill=ln-tip!25] (62,3) rectangle (88,9) node[midway] {PFN};
  \draw[fill=ln-box] (88,3) rectangle (112,9) node[midway] {\iflangja{オフセット}{offset}};
  \draw[->, ln-accent, thick] (85,25.5) -- (87,25.5) -- (87,9);
  \draw[->] (100,44) -- (100,9);
\end{tikzpicture}

% Width: body (120 mm); the drawing is about 116 mm wide.
\begin{tikzpicture}[x=0.95mm, y=1mm, font=\scriptsize\sffamily,
    >={Stealth[length=1.6mm]}]
  % virtual address fields
  \node[anchor=east] at (-1,47) {VA};
  \foreach \i/\lab/\bits in {0/L4/47--39, 1/L3/38--30, 2/L2/29--21, 3/L1/20--12} {
    \draw[fill=ln-accent!20] (\i*22,44) rectangle ++(22,6)
      node[midway] {\lab\ \iflangja{インデックス}{index}};
    \node[anchor=south, text=ln-muted] at (\i*22+11,50) {\bits};
  }
  \draw[fill=ln-box] (88,44) rectangle (112,50)
    node[midway] {\iflangja{オフセット}{offset}};
  \node[anchor=south, text=ln-muted] at (100,50) {11--0};
  % tables: three rows each, the middle one selected
  \foreach \i/\lab/\entry in {0/L4/{L3 base}, 1/L3/{L2 base}, 2/L2/{L1 base}, 3/L1/{PTE: PFN}} {
    \draw[fill=white] (\i*22+3,32) rectangle ++(16,-4) node[midway] {$\cdots$};
    \draw[fill=ln-tip!25] (\i*22+3,28) rectangle ++(16,-5) node[midway] {\entry};
    \draw[fill=white] (\i*22+3,23) rectangle ++(16,-4) node[midway] {$\cdots$};
    \node[text=ln-muted] at (\i*22+11,16.5) {\lab\ \iflangja{テーブル}{table}};
    \draw[->] (\i*22+11,44) -- (\i*22+11,32);
  }
  % next-table pointers
  \foreach \i in {0,1,2} {
    \draw[->, ln-accent, thick] (\i*22+19,25.5) -- ++(2,0) |- (\i*22+25,31);
  }
  % CR3
  \draw[->, ln-accent, thick] (-2,30) node[left, text=black] {CR3} -- (3,30);
  % physical address
  \node[anchor=east] at (61,6) {PA};
  \draw[fill=ln-tip!25] (62,3) rectangle (88,9) node[midway] {PFN};
  \draw[fill=ln-box] (88,3) rectangle (112,9) node[midway] {\iflangja{オフセット}{offset}};
  \draw[->, ln-accent, thick] (85,25.5) -- (87,25.5) -- (87,9);
  \draw[->] (100,44) -- (100,9);
\end{tikzpicture}

% Width: body (120 mm); the drawing is about 116 mm wide.
\begin{tikzpicture}[x=0.95mm, y=1mm, font=\scriptsize\sffamily,
    >={Stealth[length=1.6mm]}]
  % virtual address fields
  \node[anchor=east] at (-1,47) {VA};
  \foreach \i/\lab/\bits in {0/L4/47--39, 1/L3/38--30, 2/L2/29--21, 3/L1/20--12} {
    \draw[fill=ln-accent!20] (\i*22,44) rectangle ++(22,6)
      node[midway] {\lab\ \iflangja{インデックス}{index}};
    \node[anchor=south, text=ln-muted] at (\i*22+11,50) {\bits};
  }
  \draw[fill=ln-box] (88,44) rectangle (112,50)
    node[midway] {\iflangja{オフセット}{offset}};
  \node[anchor=south, text=ln-muted] at (100,50) {11--0};
  % tables: three rows each, the middle one selected
  \foreach \i/\lab/\entry in {0/L4/{L3 base}, 1/L3/{L2 base}, 2/L2/{L1 base}, 3/L1/{PTE: PFN}} {
    \draw[fill=white] (\i*22+3,32) rectangle ++(16,-4) node[midway] {$\cdots$};
    \draw[fill=ln-tip!25] (\i*22+3,28) rectangle ++(16,-5) node[midway] {\entry};
    \draw[fill=white] (\i*22+3,23) rectangle ++(16,-4) node[midway] {$\cdots$};
    \node[text=ln-muted] at (\i*22+11,16.5) {\lab\ \iflangja{テーブル}{table}};
    \draw[->] (\i*22+11,44) -- (\i*22+11,32);
  }
  % next-table pointers
  \foreach \i in {0,1,2} {
    \draw[->, ln-accent, thick] (\i*22+19,25.5) -- ++(2,0) |- (\i*22+25,31);
  }
  % CR3
  \draw[->, ln-accent, thick] (-2,30) node[left, text=black] {CR3} -- (3,30);
  % physical address
  \node[anchor=east] at (61,6) {PA};
  \draw[fill=ln-tip!25] (62,3) rectangle (88,9) node[midway] {PFN};
  \draw[fill=ln-box] (88,3) rectangle (112,9) node[midway] {\iflangja{オフセット}{offset}};
  \draw[->, ln-accent, thick] (85,25.5) -- (87,25.5) -- (87,9);
  \draw[->] (100,44) -- (100,9);
\end{tikzpicture}

  \caption{A four-level page table walk with the x86-64 layout.  Each
    index selects an entry holding the base of the next table; the
    level-1 entry is the PTE, whose PFN joined with the offset is the
    physical address.}
  \label{fig:l07-page-walk}
\end{figure}

\begin{example}
  \label{ex:l07-table-size}
  Take the x86-64 layout: a 36-bit VPN and \SI{8}{\byte} entries.
  \begin{itemize}
    \item \emph{One level}: $2^{36} \approx \num{6.9e10}$ entries,
      \SI{512}{\gibi\byte} of page table per process, nearly all of it
      for pages that are never used.
    \item \emph{Four levels, best case}: a program whose pages all lie
      under one level-1 table needs one table per level,
      $2^9+2^9+2^9+2^9 = \num{2048}$ entries in
      \SI{16}{\kibi\byte}, and maps up to 512 pages
      (\SI{2}{\mebi\byte}).
    \item \emph{A sparse program}: two arrays far apart in the address
      space need two chains below the level-4 table, $1+2\cdot3 = 7$
      tables, still only \SI{28}{\kibi\byte}.
  \end{itemize}
  Thousands of entries replace tens of billions.\caution{The size of a
single-level table is $2^{\text{VPN bits}}$, not
$2^{\text{address bits}}$: the offset bits do not index the table.}
\end{example}

\subsection{Address translation on a GPU}

Early GPUs had no virtual memory, or only a fixed one-to-one mapping.
Virtual memory became necessary as GPUs became programmable and began
to share data with the CPU, most visibly with the unified virtual
addressing of Lesson~5.\margin{Power et al.\ point out that an AMD
A10 APU with 400 lanes \enquote{can require up to 400 unique
translations in a single cycle} \cite{power2014x86}.}  Translation works
as on a CPU, with one difference: it happens after
\emph{address coalescing}.

\begin{figure}[tb]
  \centering
  % Address translation in the GPU pipeline (Lesson 7).  Top: the
% pipeline stages of Lesson 4.  Bottom: the memory stage expanded; the
% coalescer merges the warp's requests, the L1 data cache is accessed
% with virtual addresses while the TLB translates, a TLB miss goes to
% the page table walker (or a software miss handler), and the L2 cache
% is accessed with physical addresses.
% Redraws the outline figure m7l1_gpu_vm.
% Usage: % Address translation in the GPU pipeline (Lesson 7).  Top: the
% pipeline stages of Lesson 4.  Bottom: the memory stage expanded; the
% coalescer merges the warp's requests, the L1 data cache is accessed
% with virtual addresses while the TLB translates, a TLB miss goes to
% the page table walker (or a software miss handler), and the L2 cache
% is accessed with physical addresses.
% Redraws the outline figure m7l1_gpu_vm.
% Usage: % Address translation in the GPU pipeline (Lesson 7).  Top: the
% pipeline stages of Lesson 4.  Bottom: the memory stage expanded; the
% coalescer merges the warp's requests, the L1 data cache is accessed
% with virtual addresses while the TLB translates, a TLB miss goes to
% the page table walker (or a software miss handler), and the L2 cache
% is accessed with physical addresses.
% Redraws the outline figure m7l1_gpu_vm.
% Usage: \input{figures/l07-gpu-mmu}
% Width: body (120 mm); the drawing is about 114 mm wide.
\begin{tikzpicture}[x=1mm, y=1mm, font=\scriptsize\sffamily,
    >={Stealth[length=1.6mm]},
    st/.style={draw, fill=ln-box, minimum width=15mm, minimum height=8mm,
               align=center, inner sep=1pt},
    un/.style={draw, minimum width=24mm, minimum height=8mm,
               align=center, inner sep=1pt}]
  % pipeline
  \foreach \lab [count=\i from 0] in
      {\iflangja{フェッチ}{fetch}, \iflangja{デコード}{decode}, RF,
       \iflangja{スケジュール}{schedule}, \iflangja{実行\\アドレス生成}{execute\\addr.\ gen.},
       \iflangja{メモリ}{memory}, \iflangja{ライトバック}{write-back}} {
    \node[st] (s\i) at (\i*16.5+7.5,44) {\lab};
  }
  \node[st, fill=ln-accent!25] at (s5) {\iflangja{メモリ}{memory}};
  \foreach \i [evaluate=\i as \j using int(\i+1)] in {0,...,5} {
    \draw[->] (s\i) -- (s\j);
  }
  % memory stage expanded
  \draw[dashed, ln-muted] (0,36) rectangle (114,0);
  \draw[dashed, ln-muted] (s5.south west) -- (0,36);
  \draw[dashed, ln-muted] (s5.south east) -- (114,36);
  \node[un, fill=ln-accent!20] (co) at (14,18) {\iflangja{アドレス\\コアレッサ}{address\\coalescer}};
  \node[un, fill=ln-tip!20] (l1) at (48,27) {\iflangja{L1 D キャッシュ\\（仮想アドレス）}{L1 D-cache\\(virtual address)}};
  \node[un, fill=ln-source!20] (tlb) at (48,9) {L1 TLB};
  \node[un, fill=ln-caution!15] (ptw) at (95,9) {\iflangja{PTW または\\TLB ミスハンドラ}{PTW or\\TLB miss handler}};
  \node[un, fill=ln-tip!20] (l2) at (95,27) {\iflangja{L2 キャッシュ\\（物理アドレス）}{L2 cache\\(physical address)}};
  \draw[->] (co.east) -- ++(4,0) |- (l1.west);
  \draw[->] (co.east) -- ++(4,0) |- (tlb.west);
  \draw[->] (l1) -- node[above, text=ln-muted] {\iflangja{ミス}{miss}} (l2);
  \draw[->] (tlb) -- node[above, text=ln-muted] {\iflangja{ミス}{miss}} (ptw);
  \draw[->] (tlb.north east) -- node[pos=0.45, above left=-1pt, text=ln-muted] {PA} (l2.south west);
  \draw[->] (ptw) -- node[right, text=ln-muted] {\iflangja{メモリ上の PTE}{PTEs in memory}} (l2);
\end{tikzpicture}

% Width: body (120 mm); the drawing is about 114 mm wide.
\begin{tikzpicture}[x=1mm, y=1mm, font=\scriptsize\sffamily,
    >={Stealth[length=1.6mm]},
    st/.style={draw, fill=ln-box, minimum width=15mm, minimum height=8mm,
               align=center, inner sep=1pt},
    un/.style={draw, minimum width=24mm, minimum height=8mm,
               align=center, inner sep=1pt}]
  % pipeline
  \foreach \lab [count=\i from 0] in
      {\iflangja{フェッチ}{fetch}, \iflangja{デコード}{decode}, RF,
       \iflangja{スケジュール}{schedule}, \iflangja{実行\\アドレス生成}{execute\\addr.\ gen.},
       \iflangja{メモリ}{memory}, \iflangja{ライトバック}{write-back}} {
    \node[st] (s\i) at (\i*16.5+7.5,44) {\lab};
  }
  \node[st, fill=ln-accent!25] at (s5) {\iflangja{メモリ}{memory}};
  \foreach \i [evaluate=\i as \j using int(\i+1)] in {0,...,5} {
    \draw[->] (s\i) -- (s\j);
  }
  % memory stage expanded
  \draw[dashed, ln-muted] (0,36) rectangle (114,0);
  \draw[dashed, ln-muted] (s5.south west) -- (0,36);
  \draw[dashed, ln-muted] (s5.south east) -- (114,36);
  \node[un, fill=ln-accent!20] (co) at (14,18) {\iflangja{アドレス\\コアレッサ}{address\\coalescer}};
  \node[un, fill=ln-tip!20] (l1) at (48,27) {\iflangja{L1 D キャッシュ\\（仮想アドレス）}{L1 D-cache\\(virtual address)}};
  \node[un, fill=ln-source!20] (tlb) at (48,9) {L1 TLB};
  \node[un, fill=ln-caution!15] (ptw) at (95,9) {\iflangja{PTW または\\TLB ミスハンドラ}{PTW or\\TLB miss handler}};
  \node[un, fill=ln-tip!20] (l2) at (95,27) {\iflangja{L2 キャッシュ\\（物理アドレス）}{L2 cache\\(physical address)}};
  \draw[->] (co.east) -- ++(4,0) |- (l1.west);
  \draw[->] (co.east) -- ++(4,0) |- (tlb.west);
  \draw[->] (l1) -- node[above, text=ln-muted] {\iflangja{ミス}{miss}} (l2);
  \draw[->] (tlb) -- node[above, text=ln-muted] {\iflangja{ミス}{miss}} (ptw);
  \draw[->] (tlb.north east) -- node[pos=0.45, above left=-1pt, text=ln-muted] {PA} (l2.south west);
  \draw[->] (ptw) -- node[right, text=ln-muted] {\iflangja{メモリ上の PTE}{PTEs in memory}} (l2);
\end{tikzpicture}

% Width: body (120 mm); the drawing is about 114 mm wide.
\begin{tikzpicture}[x=1mm, y=1mm, font=\scriptsize\sffamily,
    >={Stealth[length=1.6mm]},
    st/.style={draw, fill=ln-box, minimum width=15mm, minimum height=8mm,
               align=center, inner sep=1pt},
    un/.style={draw, minimum width=24mm, minimum height=8mm,
               align=center, inner sep=1pt}]
  % pipeline
  \foreach \lab [count=\i from 0] in
      {\iflangja{フェッチ}{fetch}, \iflangja{デコード}{decode}, RF,
       \iflangja{スケジュール}{schedule}, \iflangja{実行\\アドレス生成}{execute\\addr.\ gen.},
       \iflangja{メモリ}{memory}, \iflangja{ライトバック}{write-back}} {
    \node[st] (s\i) at (\i*16.5+7.5,44) {\lab};
  }
  \node[st, fill=ln-accent!25] at (s5) {\iflangja{メモリ}{memory}};
  \foreach \i [evaluate=\i as \j using int(\i+1)] in {0,...,5} {
    \draw[->] (s\i) -- (s\j);
  }
  % memory stage expanded
  \draw[dashed, ln-muted] (0,36) rectangle (114,0);
  \draw[dashed, ln-muted] (s5.south west) -- (0,36);
  \draw[dashed, ln-muted] (s5.south east) -- (114,36);
  \node[un, fill=ln-accent!20] (co) at (14,18) {\iflangja{アドレス\\コアレッサ}{address\\coalescer}};
  \node[un, fill=ln-tip!20] (l1) at (48,27) {\iflangja{L1 D キャッシュ\\（仮想アドレス）}{L1 D-cache\\(virtual address)}};
  \node[un, fill=ln-source!20] (tlb) at (48,9) {L1 TLB};
  \node[un, fill=ln-caution!15] (ptw) at (95,9) {\iflangja{PTW または\\TLB ミスハンドラ}{PTW or\\TLB miss handler}};
  \node[un, fill=ln-tip!20] (l2) at (95,27) {\iflangja{L2 キャッシュ\\（物理アドレス）}{L2 cache\\(physical address)}};
  \draw[->] (co.east) -- ++(4,0) |- (l1.west);
  \draw[->] (co.east) -- ++(4,0) |- (tlb.west);
  \draw[->] (l1) -- node[above, text=ln-muted] {\iflangja{ミス}{miss}} (l2);
  \draw[->] (tlb) -- node[above, text=ln-muted] {\iflangja{ミス}{miss}} (ptw);
  \draw[->] (tlb.north east) -- node[pos=0.45, above left=-1pt, text=ln-muted] {PA} (l2.south west);
  \draw[->] (ptw) -- node[right, text=ln-muted] {\iflangja{メモリ上の PTE}{PTEs in memory}} (l2);
\end{tikzpicture}

  \caption{The memory stage of the GPU pipeline with address
    translation.  The coalescer merges the warp's requests; the L1
    data cache is accessed with virtual addresses while the TLB
    translates; a TLB miss goes to the PTW or a miss handler; the L2
    cache is accessed with physical addresses.}
  \label{fig:l07-gpu-mmu}
\end{figure}

In the pipeline of Lesson~4 the memory stage now contains the steps of
\cref{fig:l07-gpu-mmu}.  The coalescer merges the requests of the
warp's threads (Lesson~4); each resulting request needs a translation,
looked up in the TLB; a miss is resolved by the PTW or by a TLB miss
handler.  In this organisation the L1 data cache is accessed with
virtual addresses and the L2 cache, shared by all SMs, with physical
addresses.\margin{Both placements exist: the proof-of-concept MMU of
Power et al.\ translates after the coalescer and \emph{before} the L1,
so its L1 is physically addressed \cite{power2014x86}.}

Coalescing decides how many translations a warp instruction needs.
Two pages that hold requests of the same instruction need two TLB
lookups, however close the requests are; two requests in one page need
one.

\begin{example}
  \label{ex:l07-tlb-coalescing}
  Take a warp of eight threads, \SI{64}{\byte} memory requests and a
  toy page size of \SI{256}{\byte}, so that the VPN is the address
  shifted right by 8 bits.\margin{The toy page of \SI{256}{\byte}
keeps the numbers small; with \SI{4}{\kibi\byte} pages the same effect
appears with strides of kilobytes, such as the column of a row-major
matrix with \SI{4}{\kibi\byte} rows.}  A load through \texttt{R1} reads
  the consecutive words 64, 68, \ldots, 92: one request, one page, one
  TLB lookup.  A load through \texttt{R2} reads with a stride of
  \SI{160}{\byte}:
  \begin{center}
    \small
    \begin{tabular}{@{}lcccccccc@{}}
      \toprule
      thread & 0 & 1 & 2 & 3 & 4 & 5 & 6 & 7 \\
      \midrule
      address & 0 & 160 & 320 & 480 & 640 & 800 & 960 & 1120 \\
      hex & \texttt{000} & \texttt{0A0} & \texttt{140} & \texttt{1E0} & \texttt{280} & \texttt{320} & \texttt{3C0} & \texttt{460} \\
      VPN & 0 & 0 & 1 & 1 & 2 & 3 & 3 & 4 \\
      \bottomrule
    \end{tabular}
  \end{center}
  No two addresses fall into one \SI{64}{\byte} request, so the load
  issues eight requests; they touch five distinct pages, so after
  merging requests to the same page it needs five TLB lookups for a
  single instruction.  With a stride of \SI{256}{\byte} or more every
  thread would touch its own page.
\end{example}

In the worst case an uncoalesced load needs as many translations as
the warp has threads.\tip{TLB lookups per warp instruction: the
number of distinct pages its addresses touch, at most the number of
threads in the warp; a coalesced access needs one or two.}  Uncoalesced accesses
therefore cost twice: once in memory bandwidth (Lesson~4) and once in
translation bandwidth.

\subsection{Scalability of translation}

The challenge of GPU address translation is scalability.  One warp
instruction can produce many TLB misses at once, and dozens of warps
per SM and dozens of SMs issue them concurrently; the page table walker
must therefore sustain a high bandwidth of walks.\margin{Measured on
the Rodinia benchmarks and a database sort: coalescing and the scratchpad cut the translation
rate by 6.8$\times$, TLB misses come in bursts of 60 concurrent misses
on average, and post-coalescer TLBs miss 29\,\% of the time
\cite{power2014x86}.}  \Cref{fig:l07-ptw-sharing} shows two
organisations.  With a private PTW per SM each SM walks its own
misses, but every SM pays for a walker that idles most of the time.
With a shared PTW the SMs keep only their L1 TLBs and send their misses
to one high-bandwidth walker next to a shared L2 TLB.

\begin{figure}[tb]
  \centering
  % Private versus shared page table walkers (Lesson 7).  (a) Every SM
% has an L1 TLB and its own PTW; an L2 TLB is shared.  (b) SMs keep
% only an L1 TLB; one high-bandwidth PTW (with a page walk cache) sits
% beside the shared L2 TLB: the lecture's organisation.  The design of
% Power et al. (2014) has the shared PTW and PWC but no L2 TLB.
% Redraws the outline figures m7l1_challenges and m7l1_solutions.
% Usage: % Private versus shared page table walkers (Lesson 7).  (a) Every SM
% has an L1 TLB and its own PTW; an L2 TLB is shared.  (b) SMs keep
% only an L1 TLB; one high-bandwidth PTW (with a page walk cache) sits
% beside the shared L2 TLB: the lecture's organisation.  The design of
% Power et al. (2014) has the shared PTW and PWC but no L2 TLB.
% Redraws the outline figures m7l1_challenges and m7l1_solutions.
% Usage: % Private versus shared page table walkers (Lesson 7).  (a) Every SM
% has an L1 TLB and its own PTW; an L2 TLB is shared.  (b) SMs keep
% only an L1 TLB; one high-bandwidth PTW (with a page walk cache) sits
% beside the shared L2 TLB: the lecture's organisation.  The design of
% Power et al. (2014) has the shared PTW and PWC but no L2 TLB.
% Redraws the outline figures m7l1_challenges and m7l1_solutions.
% Usage: \input{figures/l07-ptw-sharing}
% Width: body (120 mm); the drawing is about 116 mm wide.
\begin{tikzpicture}[x=1mm, y=1mm, font=\scriptsize\sffamily,
    >={Stealth[length=1.6mm]},
    sm/.style={draw, fill=ln-box, minimum width=16mm, minimum height=12mm,
               align=center, inner sep=1pt},
    bx/.style={draw, minimum height=6mm, align=center, inner sep=1pt}]
  \foreach \ox/\title/\priv in {0/{\iflangja{(a) SM ごとの PTW}{(a) private PTW}}/1,
                                62/{\iflangja{(b) 共有 PTW}{(b) shared PTW}}/0} {
    \node[anchor=west, font=\scriptsize\sffamily\bfseries] at (\ox,38) {\title};
    \foreach \k in {0,1,2} {
      \node[sm] (sm\ox\k) at (\ox+9+\k*19,27) {};
      \node[anchor=north, inner sep=1pt] at (sm\ox\k.north) {SM\,\k};
      \ifnum\priv=1
        \node[bx, fill=ln-source!20, minimum width=7mm] at (\ox+5.2+\k*19,24.5) {\iflangja{TLB}{TLB}};
        \node[bx, fill=ln-caution!15, minimum width=7mm] at (\ox+12.8+\k*19,24.5) {PTW};
      \else
        \node[bx, fill=ln-source!20, minimum width=12mm] at (\ox+9+\k*19,24.5) {L1 TLB};
      \fi
      \draw[->] (sm\ox\k.south) -- (\ox+9+\k*19,15);
    }
    \node[bx, fill=ln-source!20, minimum width=54mm] (l2\ox) at (\ox+28,12) {\iflangja{共有 L2 TLB}{shared L2 TLB}};
    \ifnum\priv=1
      \node[bx, fill=white, minimum width=54mm] (mem\ox) at (\ox+28,0) {\iflangja{L2 キャッシュ／メモリ（ページテーブル）}{L2 cache / memory (page tables)}};
      \draw[->] (l2\ox.south) -- (l2\ox.south |- mem\ox.north);
    \else
      \node[bx, fill=white, minimum width=30mm, minimum height=8mm] (mem\ox) at (\ox+15,-1)
        {\iflangja{L2 キャッシュ／メモリ\\（ページテーブル）}{L2 cache / memory\\(page tables)}};
      \node[bx, fill=ln-caution!15, minimum width=17mm, minimum height=8mm] (pw) at (\ox+46.5,-1)
        {\iflangja{共有 PTW\\＋PWC}{shared PTW\\+ PWC}};
      \draw[->] (pw.north |- l2\ox.south) -- node[right, text=ln-muted] {\iflangja{ミス}{miss}} (pw.north);
      \draw[<->] (pw.west) -- (mem\ox.east);
    \fi
  }
\end{tikzpicture}

% Width: body (120 mm); the drawing is about 116 mm wide.
\begin{tikzpicture}[x=1mm, y=1mm, font=\scriptsize\sffamily,
    >={Stealth[length=1.6mm]},
    sm/.style={draw, fill=ln-box, minimum width=16mm, minimum height=12mm,
               align=center, inner sep=1pt},
    bx/.style={draw, minimum height=6mm, align=center, inner sep=1pt}]
  \foreach \ox/\title/\priv in {0/{\iflangja{(a) SM ごとの PTW}{(a) private PTW}}/1,
                                62/{\iflangja{(b) 共有 PTW}{(b) shared PTW}}/0} {
    \node[anchor=west, font=\scriptsize\sffamily\bfseries] at (\ox,38) {\title};
    \foreach \k in {0,1,2} {
      \node[sm] (sm\ox\k) at (\ox+9+\k*19,27) {};
      \node[anchor=north, inner sep=1pt] at (sm\ox\k.north) {SM\,\k};
      \ifnum\priv=1
        \node[bx, fill=ln-source!20, minimum width=7mm] at (\ox+5.2+\k*19,24.5) {\iflangja{TLB}{TLB}};
        \node[bx, fill=ln-caution!15, minimum width=7mm] at (\ox+12.8+\k*19,24.5) {PTW};
      \else
        \node[bx, fill=ln-source!20, minimum width=12mm] at (\ox+9+\k*19,24.5) {L1 TLB};
      \fi
      \draw[->] (sm\ox\k.south) -- (\ox+9+\k*19,15);
    }
    \node[bx, fill=ln-source!20, minimum width=54mm] (l2\ox) at (\ox+28,12) {\iflangja{共有 L2 TLB}{shared L2 TLB}};
    \ifnum\priv=1
      \node[bx, fill=white, minimum width=54mm] (mem\ox) at (\ox+28,0) {\iflangja{L2 キャッシュ／メモリ（ページテーブル）}{L2 cache / memory (page tables)}};
      \draw[->] (l2\ox.south) -- (l2\ox.south |- mem\ox.north);
    \else
      \node[bx, fill=white, minimum width=30mm, minimum height=8mm] (mem\ox) at (\ox+15,-1)
        {\iflangja{L2 キャッシュ／メモリ\\（ページテーブル）}{L2 cache / memory\\(page tables)}};
      \node[bx, fill=ln-caution!15, minimum width=17mm, minimum height=8mm] (pw) at (\ox+46.5,-1)
        {\iflangja{共有 PTW\\＋PWC}{shared PTW\\+ PWC}};
      \draw[->] (pw.north |- l2\ox.south) -- node[right, text=ln-muted] {\iflangja{ミス}{miss}} (pw.north);
      \draw[<->] (pw.west) -- (mem\ox.east);
    \fi
  }
\end{tikzpicture}

% Width: body (120 mm); the drawing is about 116 mm wide.
\begin{tikzpicture}[x=1mm, y=1mm, font=\scriptsize\sffamily,
    >={Stealth[length=1.6mm]},
    sm/.style={draw, fill=ln-box, minimum width=16mm, minimum height=12mm,
               align=center, inner sep=1pt},
    bx/.style={draw, minimum height=6mm, align=center, inner sep=1pt}]
  \foreach \ox/\title/\priv in {0/{\iflangja{(a) SM ごとの PTW}{(a) private PTW}}/1,
                                62/{\iflangja{(b) 共有 PTW}{(b) shared PTW}}/0} {
    \node[anchor=west, font=\scriptsize\sffamily\bfseries] at (\ox,38) {\title};
    \foreach \k in {0,1,2} {
      \node[sm] (sm\ox\k) at (\ox+9+\k*19,27) {};
      \node[anchor=north, inner sep=1pt] at (sm\ox\k.north) {SM\,\k};
      \ifnum\priv=1
        \node[bx, fill=ln-source!20, minimum width=7mm] at (\ox+5.2+\k*19,24.5) {\iflangja{TLB}{TLB}};
        \node[bx, fill=ln-caution!15, minimum width=7mm] at (\ox+12.8+\k*19,24.5) {PTW};
      \else
        \node[bx, fill=ln-source!20, minimum width=12mm] at (\ox+9+\k*19,24.5) {L1 TLB};
      \fi
      \draw[->] (sm\ox\k.south) -- (\ox+9+\k*19,15);
    }
    \node[bx, fill=ln-source!20, minimum width=54mm] (l2\ox) at (\ox+28,12) {\iflangja{共有 L2 TLB}{shared L2 TLB}};
    \ifnum\priv=1
      \node[bx, fill=white, minimum width=54mm] (mem\ox) at (\ox+28,0) {\iflangja{L2 キャッシュ／メモリ（ページテーブル）}{L2 cache / memory (page tables)}};
      \draw[->] (l2\ox.south) -- (l2\ox.south |- mem\ox.north);
    \else
      \node[bx, fill=white, minimum width=30mm, minimum height=8mm] (mem\ox) at (\ox+15,-1)
        {\iflangja{L2 キャッシュ／メモリ\\（ページテーブル）}{L2 cache / memory\\(page tables)}};
      \node[bx, fill=ln-caution!15, minimum width=17mm, minimum height=8mm] (pw) at (\ox+46.5,-1)
        {\iflangja{共有 PTW\\＋PWC}{shared PTW\\+ PWC}};
      \draw[->] (pw.north |- l2\ox.south) -- node[right, text=ln-muted] {\iflangja{ミス}{miss}} (pw.north);
      \draw[<->] (pw.west) -- (mem\ox.east);
    \fi
  }
\end{tikzpicture}

  \caption{(a) A private PTW in every SM, with a shared L2 TLB.
    (b) SMs with L1 TLBs only and one shared, high-bandwidth PTW with
    a page walk cache (PWC) next to the shared L2 TLB.  The design of
    Power et al.\ is (b) without the L2 TLB.}
  \label{fig:l07-ptw-sharing}
\end{figure}

Power, Hill and Wood build the shared organisation from known CPU
parts \cite{power2014x86}: per-compute-unit TLBs placed after the
coalescer, a highly threaded PTW that keeps the state of many
outstanding walks in a set of buffers, and a
\term{page walk cache}[ページウォークキャッシュ], a small cache in the
MMU that holds entries of the upper (non-leaf) levels of the page
table, so that most walks need fewer than four memory
accesses.\margin{The page walk cache cuts the average walk latency by
over 95\,\%; the walker needs it because the GPU's L2, where the PTEs
are cached, takes nearly 300 cycles.  The complete design comes within
2\,\% of an ideal MMU on average, with \SI{4}{\kibi\byte} pages and
unmodified x86-64 page tables \cite{power2014x86}.}  Their design has
no L2 TLB: the per-CU TLBs feed the shared walker and its page walk
cache directly.  A shared L2 TLB in place of the page walk cache was
more than 2$\times$ slower on average, and one added beside it changed
performance by about 0.1\,\% \cite[§6]{power2014x86}.

\subsection{Reducing the cost of translation}

Three techniques reduce the number or the cost of
translations.
\begin{marginfigure}
  % Segment mapping (Lesson 7): three programs, each with the virtual
% space 0x0000--0xFFFF, occupy three contiguous 64 KiB segments of
% physical memory; translation adds the segment base.
% Redraws the outline figure m7l1_solutions (left half).
% Usage: % Segment mapping (Lesson 7): three programs, each with the virtual
% space 0x0000--0xFFFF, occupy three contiguous 64 KiB segments of
% physical memory; translation adds the segment base.
% Redraws the outline figure m7l1_solutions (left half).
% Usage: % Segment mapping (Lesson 7): three programs, each with the virtual
% space 0x0000--0xFFFF, occupy three contiguous 64 KiB segments of
% physical memory; translation adds the segment base.
% Redraws the outline figure m7l1_solutions (left half).
% Usage: \input{figures/l07-segments}  (in a marginfigure)
% Width: margin (50 mm); the drawing is about 48 mm wide.
\begin{tikzpicture}[x=1mm, y=1mm, font=\scriptsize\sffamily]
  \node[anchor=south, text=ln-muted] at (8,30) {\iflangja{物理アドレス}{physical}};
  \foreach \k/\col/\lo/\hi/\prog in {0/ln-caution!20/000000/00FFFF/3,
                                     1/ln-accent!20/010000/01FFFF/2,
                                     2/ln-tip!25/020000/02FFFF/1} {
    \draw[fill=\col] (0,30-\k*10) rectangle ++(16,-10);
    \node[anchor=north west, inner sep=1pt, font=\scriptsize\ttfamily] at (0.3,30-\k*10) {\lo};
    \node[anchor=south west, inner sep=1pt, font=\scriptsize\ttfamily] at (0.3,20-\k*10) {\hi};
    \node[anchor=west, align=left] at (18,25-\k*10)
      {\iflangja{プログラム}{program}~\prog\\[-1pt]VA {\ttfamily 0000}--{\ttfamily FFFF}};
  }
\end{tikzpicture}
  (in a marginfigure)
% Width: margin (50 mm); the drawing is about 48 mm wide.
\begin{tikzpicture}[x=1mm, y=1mm, font=\scriptsize\sffamily]
  \node[anchor=south, text=ln-muted] at (8,30) {\iflangja{物理アドレス}{physical}};
  \foreach \k/\col/\lo/\hi/\prog in {0/ln-caution!20/000000/00FFFF/3,
                                     1/ln-accent!20/010000/01FFFF/2,
                                     2/ln-tip!25/020000/02FFFF/1} {
    \draw[fill=\col] (0,30-\k*10) rectangle ++(16,-10);
    \node[anchor=north west, inner sep=1pt, font=\scriptsize\ttfamily] at (0.3,30-\k*10) {\lo};
    \node[anchor=south west, inner sep=1pt, font=\scriptsize\ttfamily] at (0.3,20-\k*10) {\hi};
    \node[anchor=west, align=left] at (18,25-\k*10)
      {\iflangja{プログラム}{program}~\prog\\[-1pt]VA {\ttfamily 0000}--{\ttfamily FFFF}};
  }
\end{tikzpicture}
  (in a marginfigure)
% Width: margin (50 mm); the drawing is about 48 mm wide.
\begin{tikzpicture}[x=1mm, y=1mm, font=\scriptsize\sffamily]
  \node[anchor=south, text=ln-muted] at (8,30) {\iflangja{物理アドレス}{physical}};
  \foreach \k/\col/\lo/\hi/\prog in {0/ln-caution!20/000000/00FFFF/3,
                                     1/ln-accent!20/010000/01FFFF/2,
                                     2/ln-tip!25/020000/02FFFF/1} {
    \draw[fill=\col] (0,30-\k*10) rectangle ++(16,-10);
    \node[anchor=north west, inner sep=1pt, font=\scriptsize\ttfamily] at (0.3,30-\k*10) {\lo};
    \node[anchor=south west, inner sep=1pt, font=\scriptsize\ttfamily] at (0.3,20-\k*10) {\hi};
    \node[anchor=west, align=left] at (18,25-\k*10)
      {\iflangja{プログラム}{program}~\prog\\[-1pt]VA {\ttfamily 0000}--{\ttfamily FFFF}};
  }
\end{tikzpicture}

  \caption{Segment mapping: three programs with the same virtual range
    in three \SI{64}{\kibi\byte} physical segments.}
  \label{fig:l07-segments}
\end{marginfigure}
\begin{description}
  \item[Segment mapping.]  In \term{segment mapping}[セグメントマッピング]
    each program's virtual space maps onto one contiguous region of
    physical memory, so translation is the addition of a base
    (\cref{fig:l07-segments}): in the figure, virtual address
    \texttt{0x1234} of program~2 is physical address
    \texttt{0x011234}.  No table and no walk are needed, but each
    program's space must be contiguous in physical memory.
  \item[Large pages.]  Fewer, larger pages need fewer translations
    (below).
  \item[Multi-level TLBs.]  A small L1 TLB per SM backed by a larger,
    shared L2 TLB, as in \cref{fig:l07-ptw-sharing}, can serve an L1
    miss with a translation that another SM already brought in, before
    a walk starts (with a page walk cache present, Power et al.\ saw no
    gain from it, above).
\end{description}

\begin{definition}[Large page]
  \label{def:l07-large-page}
  A \term{large page}[ラージページ] is a page much larger than the
  base page, typically \SI{2}{\mebi\byte} instead of
  \SI{4}{\kibi\byte}.  Its offset is 21 bits instead of 12, so the VPN
  is 9 bits shorter, and on x86-64 the walk ends one level early.  One
  TLB entry then covers $\SI{2}{\mebi\byte}/\SI{4}{\kibi\byte} = 512$
  times as much memory.
\end{definition}

The amount of memory a TLB can map at once, entries times page size,
is its \emph{reach}.\margin{CPG §19.2.2.1 Memory Paging and Page
Sizes: x86-64 CPUs use \SI{4}{\kibi\byte} physical pages, while
\enquote{NVIDIA GPUs support multiple physical page sizes, but prefer
2MiB physical pages or larger}; TLB misses \enquote{are generally
significantly more expensive on the GPU compared to the CPU}
\cite{nvidia2024cudaguide}.}  The cost of a large page is
\term{internal fragmentation}[内部断片化]: memory allocated in whole
pages but not used, lost inside the last page of an allocation.

\begin{example}
  \label{ex:l07-large-pages}
  An L1 TLB of 32 entries reaches $32 \times \SI{4}{\kibi\byte} =
  \SI{128}{\kibi\byte}$ with base pages and $32 \times
  \SI{2}{\mebi\byte} = \SI{64}{\mebi\byte}$ with large pages.  An
  allocation of \SI{10}{\kibi\byte} occupies three \SI{4}{\kibi\byte}
  pages and wastes \SI{2}{\kibi\byte}, but one \SI{2}{\mebi\byte} page
  wastes \SI{2038}{\kibi\byte}, over 99\,\% of the page (our numbers,
  chosen for the arithmetic).
\end{example}

Large and small pages can also be mixed, so that each region of
memory gets the page size that suits it.\sidenote{Research has continued along these lines.
Mosaic \cite{ausavarungnirun2017mosaic} supports several page sizes at
once without application changes and reports 55.5\,\% (homogeneous)
and 29.7\,\% (heterogeneous) average gains for multi-application
workloads, within 6.8\,\% and 15.4\,\% of an ideal TLB.  For several
GPUs, Li et al.\ \cite{li2021tlb} share translations among peer GPUs
and spill TLB entries to the IOMMU's TLB, gaining 23.5\,\% and
16.3\,\% on average for single- and multi-application workloads.}

\begin{keypoint}
  A GPU translates the addresses of whole warps: coalescing decides how
  many translations an instruction needs, an uncoalesced access can
  need one per thread, and the walker must handle many misses at once.
  Shared high-bandwidth walkers with page walk caches, multi-level
  TLBs, segments and large pages reduce the cost; large pages trade
  TLB reach against internal fragmentation.
\end{keypoint}

\section{Unified virtual address space and on-demand paging}
\label{sec:l07-uva}

The second side of virtual memory is page allocation, and here the GPU
differs from the CPU most.\source{M7, UVA and on-demand paging;
CPG §19 Unified Memory Programming; BPG §10.1.4 Unified Virtual
Addressing; \cite{ganguly2019interplay,kim2020batch}.}  This section
compares explicit copying with a shared address space, follows a GPU
page fault to the host and back, and presents the prefetching and
eviction policies that make on-demand paging affordable.

\subsection{Copy-then-execute versus a shared address space}

\begin{definition}[Copy-then-execute model]
  \label{def:l07-copy-then-execute}
  In the \term{copy-then-execute model}[コピー後実行モデル] the host
  allocates device memory (\texttt{cudaMalloc}) and copies the input
  into it (\texttt{cudaMemcpy}) before it launches the kernel, and
  copies the results back afterwards.  All data a kernel touches is in
  device memory before the kernel starts.
\end{definition}

The data transfer is explicit and happens in the preparation before
the kernel launch.  It has one strength: bulk transfers that the
programmer can schedule use the PCIe bandwidth fully, as the pinned
buffers and streams of Lesson~5 showed.  Its weakness is that the
programmer manages every buffer and every copy by hand, and that data
structures with pointers cannot be shared at all.

\caution{\enquote{SVM} is the term of OpenCL 2.0.  CUDA does not use
it: it calls the shared address space unified virtual addressing and
the migrating variant unified (managed) memory.}%
\begin{definition}[Shared virtual memory]
  \label{def:l07-svm}
  With \term{shared virtual memory (SVM)}[共有仮想メモリ], the name
  used in the OpenCL/CUDA programming model, the CPU and the GPU use
  one virtual address space, so that one pointer means the same datum
  on both (\cref{fig:l07-uva}).\index{SVM|see{shared virtual memory (SVM)}}
\end{definition}

\begin{figure}[tb]
  \centering
  % Separate memory spaces versus a unified virtual address space
% (Lesson 7).  (a) The CPU writes and reads a buffer in CPU memory, the
% GPU one in GPU memory, and explicit copies move data between them.
% (b) One logical buffer in a unified space is read and written by
% both processors; no copy appears in the program.
% Redraws the outline figure m7l2_uva.
% Usage: % Separate memory spaces versus a unified virtual address space
% (Lesson 7).  (a) The CPU writes and reads a buffer in CPU memory, the
% GPU one in GPU memory, and explicit copies move data between them.
% (b) One logical buffer in a unified space is read and written by
% both processors; no copy appears in the program.
% Redraws the outline figure m7l2_uva.
% Usage: % Separate memory spaces versus a unified virtual address space
% (Lesson 7).  (a) The CPU writes and reads a buffer in CPU memory, the
% GPU one in GPU memory, and explicit copies move data between them.
% (b) One logical buffer in a unified space is read and written by
% both processors; no copy appears in the program.
% Redraws the outline figure m7l2_uva.
% Usage: \input{figures/l07-uva}
% Width: body (120 mm); the drawing is about 112 mm wide.
\begin{tikzpicture}[x=1mm, y=1mm, font=\scriptsize\sffamily,
    >={Stealth[length=1.6mm]},
    pr/.style={draw, fill=ln-accent!20, minimum width=12mm, minimum height=6mm},
    buf/.style={draw, fill=ln-tip!25, minimum width=14mm, minimum height=6mm}]
  % (a) separate
  \node[anchor=west, font=\scriptsize\sffamily\bfseries] at (0,33)
    {\iflangja{(a) 分離したメモリ空間}{(a) separate memory spaces}};
  \node[pr] (c1) at (9,24) {CPU};
  \node[pr] (g1) at (45,24) {GPU};
  \draw[fill=ln-box] (0,15) rectangle (20,3);
  \node[anchor=north, text=ln-muted] at (10,15) {\iflangja{CPU メモリ}{CPU memory}};
  \node[buf] (b1) at (10,7) {\iflangja{バッファ}{buffer}};
  \draw[fill=ln-box] (34,15) rectangle (54,3);
  \node[anchor=north, text=ln-muted] at (44,15) {\iflangja{GPU メモリ}{GPU memory}};
  \node[buf] (b2) at (44,7) {\iflangja{バッファ}{buffer}};
  \draw[<->] (c1) -- (c1 |- 0,15);
  \draw[<->] (g1) -- (g1 |- 0,15);
  \draw[->, thick, ln-caution] (b1.east) -- node[above, text=ln-caution] {\iflangja{コピー}{copy}} (b2.west);
  % (b) unified
  \node[anchor=west, font=\scriptsize\sffamily\bfseries] at (62,33)
    {\iflangja{(b) 統合仮想アドレス空間}{(b) unified virtual address space}};
  \node[pr] (c2) at (71,24) {CPU};
  \node[pr] (g2) at (103,24) {GPU};
  \draw[fill=ln-box] (62,15) rectangle (112,3);
  \node[anchor=north, text=ln-muted] at (87,2.5) {\iflangja{統合メモリ空間}{unified memory space}};
  \node[buf] (b3) at (87,7) {\iflangja{バッファ}{buffer}};
  \draw[<->] (c2) -- (b3);
  \draw[<->] (g2) -- (b3);
\end{tikzpicture}

% Width: body (120 mm); the drawing is about 112 mm wide.
\begin{tikzpicture}[x=1mm, y=1mm, font=\scriptsize\sffamily,
    >={Stealth[length=1.6mm]},
    pr/.style={draw, fill=ln-accent!20, minimum width=12mm, minimum height=6mm},
    buf/.style={draw, fill=ln-tip!25, minimum width=14mm, minimum height=6mm}]
  % (a) separate
  \node[anchor=west, font=\scriptsize\sffamily\bfseries] at (0,33)
    {\iflangja{(a) 分離したメモリ空間}{(a) separate memory spaces}};
  \node[pr] (c1) at (9,24) {CPU};
  \node[pr] (g1) at (45,24) {GPU};
  \draw[fill=ln-box] (0,15) rectangle (20,3);
  \node[anchor=north, text=ln-muted] at (10,15) {\iflangja{CPU メモリ}{CPU memory}};
  \node[buf] (b1) at (10,7) {\iflangja{バッファ}{buffer}};
  \draw[fill=ln-box] (34,15) rectangle (54,3);
  \node[anchor=north, text=ln-muted] at (44,15) {\iflangja{GPU メモリ}{GPU memory}};
  \node[buf] (b2) at (44,7) {\iflangja{バッファ}{buffer}};
  \draw[<->] (c1) -- (c1 |- 0,15);
  \draw[<->] (g1) -- (g1 |- 0,15);
  \draw[->, thick, ln-caution] (b1.east) -- node[above, text=ln-caution] {\iflangja{コピー}{copy}} (b2.west);
  % (b) unified
  \node[anchor=west, font=\scriptsize\sffamily\bfseries] at (62,33)
    {\iflangja{(b) 統合仮想アドレス空間}{(b) unified virtual address space}};
  \node[pr] (c2) at (71,24) {CPU};
  \node[pr] (g2) at (103,24) {GPU};
  \draw[fill=ln-box] (62,15) rectangle (112,3);
  \node[anchor=north, text=ln-muted] at (87,2.5) {\iflangja{統合メモリ空間}{unified memory space}};
  \node[buf] (b3) at (87,7) {\iflangja{バッファ}{buffer}};
  \draw[<->] (c2) -- (b3);
  \draw[<->] (g2) -- (b3);
\end{tikzpicture}

% Width: body (120 mm); the drawing is about 112 mm wide.
\begin{tikzpicture}[x=1mm, y=1mm, font=\scriptsize\sffamily,
    >={Stealth[length=1.6mm]},
    pr/.style={draw, fill=ln-accent!20, minimum width=12mm, minimum height=6mm},
    buf/.style={draw, fill=ln-tip!25, minimum width=14mm, minimum height=6mm}]
  % (a) separate
  \node[anchor=west, font=\scriptsize\sffamily\bfseries] at (0,33)
    {\iflangja{(a) 分離したメモリ空間}{(a) separate memory spaces}};
  \node[pr] (c1) at (9,24) {CPU};
  \node[pr] (g1) at (45,24) {GPU};
  \draw[fill=ln-box] (0,15) rectangle (20,3);
  \node[anchor=north, text=ln-muted] at (10,15) {\iflangja{CPU メモリ}{CPU memory}};
  \node[buf] (b1) at (10,7) {\iflangja{バッファ}{buffer}};
  \draw[fill=ln-box] (34,15) rectangle (54,3);
  \node[anchor=north, text=ln-muted] at (44,15) {\iflangja{GPU メモリ}{GPU memory}};
  \node[buf] (b2) at (44,7) {\iflangja{バッファ}{buffer}};
  \draw[<->] (c1) -- (c1 |- 0,15);
  \draw[<->] (g1) -- (g1 |- 0,15);
  \draw[->, thick, ln-caution] (b1.east) -- node[above, text=ln-caution] {\iflangja{コピー}{copy}} (b2.west);
  % (b) unified
  \node[anchor=west, font=\scriptsize\sffamily\bfseries] at (62,33)
    {\iflangja{(b) 統合仮想アドレス空間}{(b) unified virtual address space}};
  \node[pr] (c2) at (71,24) {CPU};
  \node[pr] (g2) at (103,24) {GPU};
  \draw[fill=ln-box] (62,15) rectangle (112,3);
  \node[anchor=north, text=ln-muted] at (87,2.5) {\iflangja{統合メモリ空間}{unified memory space}};
  \node[buf] (b3) at (87,7) {\iflangja{バッファ}{buffer}};
  \draw[<->] (c2) -- (b3);
  \draw[<->] (g2) -- (b3);
\end{tikzpicture}

  \caption{(a) Separate memory spaces: each processor reads and writes
    its own buffer and explicit copies move the data.  (b) A unified
    virtual address space: one logical buffer that both processors read
    and write.}
  \label{fig:l07-uva}
\end{figure}

The \term*{unified virtual addressing (UVA)} of Lesson~5 is one address
space over separate memories; a shared address space raises
programmability further when the pages also move by
themselves.\margin{Lesson~3 introduced the programming side of
\texttt{cudaMallocManaged}; Lesson~5 separated UVA (one address space)
from unified memory (automatic page migration).}  Most GPUs in
high-performance computing are discrete: they have their own physical
memory, so a shared address space still requires copies between host
and device memory, now performed on demand.  That needs
\emph{dynamic page allocation}: the decision where a page lives is taken
while the kernel runs.

\begin{definition}[On-demand paging]
  \label{def:l07-on-demand}
  Under \term{on-demand paging}[オンデマンドページング] memory
  allocated with \texttt{cudaMallocManaged} is not placed in device
  memory in advance; the first GPU access to a page that is not
  resident faults, and the page is then allocated in device memory and
  its data migrated from the host.
\end{definition}

The two models side by side, for a kernel \texttt{scale} that
multiplies $n$ floats in place:

\begin{lstlisting}[language=C++, numbers=none]
// copy-then-execute
float *d;
cudaMalloc(&d, n * sizeof(float));
cudaMemcpy(d, h, n * sizeof(float),
           cudaMemcpyHostToDevice);
scale<<<(n + 255) / 256, 256>>>(d, n);
cudaMemcpy(h, d, n * sizeof(float),
           cudaMemcpyDeviceToHost);

// on-demand paging
float *m;
cudaMallocManaged(&m, n * sizeof(float));
for (int i = 0; i < n; ++i) m[i] = i;  // CPU touches
scale<<<(n + 255) / 256, 256>>>(m, n);  // GPU faults
cudaDeviceSynchronize();      // m is valid on the CPU
\end{lstlisting}

\margin{CPG §19.3.2.1 Data Migration and Coherency: before compute
capability 6.x, all managed memory was moved to the GPU at every kernel
launch; compute capability 6.x (Pascal) introduced GPU page faulting,
which migrates pages on demand \cite{nvidia2024cudaguide}.}%
The second version has no copy in the program, but the copies still
happen: each page crosses PCIe when the kernel first touches it.

\subsection{The page-fault path}

\begin{definition}[Far fault and GMMU]
  \label{def:l07-far-fault}
  In a GPU with unified memory, the
  \term{GPU memory management unit (GMMU)}[GPUメモリ管理ユニット]
  receives the TLB misses of the SMs and walks the GPU page table.  A
  \term{far fault}[ファーフォールト] is a GPU page fault on a page that
  is not present in device memory, which the host driver must resolve
  by allocating a page and migrating its data, in contrast to a
  \emph{near fault}, an L2 miss that device memory serves.\index{GMMU|see{GPU memory management unit (GMMU)}}
\end{definition}

\Cref{fig:l07-fault-path} follows a far fault.
\begin{enumerate}
  \item The GMMU's walk finds no valid PTE; the fault is recorded in
    the GPU's page-fault MSHRs (miss status holding registers), and
    the faulting warp waits.
  \item The fault is sent over PCIe to the CPU.
  \item The operating system and the GPU driver handle it: they
    allocate a page in device memory and migrate the data.
  \item The new mapping is sent back; the PTE is validated, and the
    access is replayed.
\end{enumerate}
The round trip takes several tens of
microseconds.\caution{Lecture vs paper: one passage of the source puts
this latency at \enquote{tens of milliseconds}; it is tens of
microseconds.  Prior work reports \SIrange{20}{50}{\micro\second} per fault, and Kim et al.\
measure \SIrange{223}{553}{\micro\second} for a whole batch of faults (median \SI{313}{\micro\second})
\cite{kim2020batch}.}

\begin{figure}[b]
  \centering
  % Far-fault handling with on-demand paging on a discrete GPU
% (Lesson 7).  (1) A TLB miss goes to the GMMU, whose walk finds no
% valid PTE: a far fault, recorded in the far-fault MSHRs.  (2) The
% fault is sent over PCIe to the host.  (3) The GPU driver and the OS
% allocate a page and migrate the data.  (4) The new mapping returns,
% the PTE is validated and the access is replayed.
% Redraws the outline figures m7l2_page_fault_handling and
% m7l2_challenge_ondemand.
% Usage: % Far-fault handling with on-demand paging on a discrete GPU
% (Lesson 7).  (1) A TLB miss goes to the GMMU, whose walk finds no
% valid PTE: a far fault, recorded in the far-fault MSHRs.  (2) The
% fault is sent over PCIe to the host.  (3) The GPU driver and the OS
% allocate a page and migrate the data.  (4) The new mapping returns,
% the PTE is validated and the access is replayed.
% Redraws the outline figures m7l2_page_fault_handling and
% m7l2_challenge_ondemand.
% Usage: % Far-fault handling with on-demand paging on a discrete GPU
% (Lesson 7).  (1) A TLB miss goes to the GMMU, whose walk finds no
% valid PTE: a far fault, recorded in the far-fault MSHRs.  (2) The
% fault is sent over PCIe to the host.  (3) The GPU driver and the OS
% allocate a page and migrate the data.  (4) The new mapping returns,
% the PTE is validated and the access is replayed.
% Redraws the outline figures m7l2_page_fault_handling and
% m7l2_challenge_ondemand.
% Usage: \input{figures/l07-fault-path}
% Width: body (120 mm); the drawing is about 116 mm wide.
\begin{tikzpicture}[x=1mm, y=1mm, font=\scriptsize\sffamily,
    >={Stealth[length=1.6mm]},
    bx/.style={draw, minimum height=6mm, align=center, inner sep=1pt},
    stp/.style={circle, draw, fill=white, inner sep=0.6pt, font=\scriptsize\sffamily\bfseries}]
  % GPU
  \draw[fill=ln-box] (0,0) rectangle (56,40);
  \node[anchor=north west, font=\scriptsize\sffamily\bfseries] at (0.5,39.5) {GPU};
  \foreach \k in {0,1,2} {
    \node[bx, fill=ln-accent!20, minimum width=14mm] at (8+\k*17,31) {SM\,\k\ +\,TLB};
  }
  \node[bx, fill=ln-caution!15, minimum width=50mm, minimum height=13mm] (gmmu) at (28,17)
    {};
  \node[anchor=north, inner sep=1pt] at (gmmu.north) {GMMU};
  \node[bx, fill=white, minimum width=14mm] at (11,15) {L2 TLB};
  \node[bx, fill=white, minimum width=14mm] at (28,15) {\iflangja{ウォーク}{walk}};
  \node[bx, fill=white, minimum width=14mm] at (45,15) {\iflangja{フォールト\\MSHR}{fault\\MSHRs}};
  \node[bx, fill=ln-tip!25, minimum width=50mm] (dmem) at (28,4.5) {\iflangja{デバイスメモリ（ページテーブル）}{device memory (page table)}};
  \foreach \k in {0,1,2} { \draw[->] (8+\k*17,28) -- (8+\k*17,23.5); }
  % PCIe
  \draw[fill=ln-muted!15] (60,14) rectangle (72,26);
  \node[rotate=90] at (66,20) {PCIe};
  % CPU
  \draw[fill=ln-box] (76,0) rectangle (116,40);
  \node[anchor=north west, font=\scriptsize\sffamily\bfseries] at (76.5,39.5) {CPU};
  \node[bx, fill=ln-accent!20, minimum width=34mm] (drv) at (96,26) {\iflangja{OS ＋ GPU ドライバ／\\ランタイム}{OS + GPU driver\\and runtime}};
  \node[bx, fill=ln-tip!25, minimum width=34mm] (hmem) at (96,7) {\iflangja{ホストメモリ}{host memory}};
  \draw[<->] (drv) -- (hmem);
  % fault and reply
  \draw[->, thick, ln-caution] (53,24) -- (80,29);
  \node[stp] at (49,25) {1};
  \node[stp] at (66,30) {2};
  \node[anchor=south, text=ln-caution, align=center] at (66,31.8) {\iflangja{ページフォールト}{page fault}};
  \node[stp] at (106,17) {3};
  \draw[->, thick, ln-tip] (80,22) -- (53,11);
  \node[stp] at (66,10) {4};
  \node[anchor=north, text=ln-tip, align=center] at (66,8.5) {\iflangja{新しいマッピング\\とデータ}{new mapping\\and data}};
\end{tikzpicture}

% Width: body (120 mm); the drawing is about 116 mm wide.
\begin{tikzpicture}[x=1mm, y=1mm, font=\scriptsize\sffamily,
    >={Stealth[length=1.6mm]},
    bx/.style={draw, minimum height=6mm, align=center, inner sep=1pt},
    stp/.style={circle, draw, fill=white, inner sep=0.6pt, font=\scriptsize\sffamily\bfseries}]
  % GPU
  \draw[fill=ln-box] (0,0) rectangle (56,40);
  \node[anchor=north west, font=\scriptsize\sffamily\bfseries] at (0.5,39.5) {GPU};
  \foreach \k in {0,1,2} {
    \node[bx, fill=ln-accent!20, minimum width=14mm] at (8+\k*17,31) {SM\,\k\ +\,TLB};
  }
  \node[bx, fill=ln-caution!15, minimum width=50mm, minimum height=13mm] (gmmu) at (28,17)
    {};
  \node[anchor=north, inner sep=1pt] at (gmmu.north) {GMMU};
  \node[bx, fill=white, minimum width=14mm] at (11,15) {L2 TLB};
  \node[bx, fill=white, minimum width=14mm] at (28,15) {\iflangja{ウォーク}{walk}};
  \node[bx, fill=white, minimum width=14mm] at (45,15) {\iflangja{フォールト\\MSHR}{fault\\MSHRs}};
  \node[bx, fill=ln-tip!25, minimum width=50mm] (dmem) at (28,4.5) {\iflangja{デバイスメモリ（ページテーブル）}{device memory (page table)}};
  \foreach \k in {0,1,2} { \draw[->] (8+\k*17,28) -- (8+\k*17,23.5); }
  % PCIe
  \draw[fill=ln-muted!15] (60,14) rectangle (72,26);
  \node[rotate=90] at (66,20) {PCIe};
  % CPU
  \draw[fill=ln-box] (76,0) rectangle (116,40);
  \node[anchor=north west, font=\scriptsize\sffamily\bfseries] at (76.5,39.5) {CPU};
  \node[bx, fill=ln-accent!20, minimum width=34mm] (drv) at (96,26) {\iflangja{OS ＋ GPU ドライバ／\\ランタイム}{OS + GPU driver\\and runtime}};
  \node[bx, fill=ln-tip!25, minimum width=34mm] (hmem) at (96,7) {\iflangja{ホストメモリ}{host memory}};
  \draw[<->] (drv) -- (hmem);
  % fault and reply
  \draw[->, thick, ln-caution] (53,24) -- (80,29);
  \node[stp] at (49,25) {1};
  \node[stp] at (66,30) {2};
  \node[anchor=south, text=ln-caution, align=center] at (66,31.8) {\iflangja{ページフォールト}{page fault}};
  \node[stp] at (106,17) {3};
  \draw[->, thick, ln-tip] (80,22) -- (53,11);
  \node[stp] at (66,10) {4};
  \node[anchor=north, text=ln-tip, align=center] at (66,8.5) {\iflangja{新しいマッピング\\とデータ}{new mapping\\and data}};
\end{tikzpicture}

% Width: body (120 mm); the drawing is about 116 mm wide.
\begin{tikzpicture}[x=1mm, y=1mm, font=\scriptsize\sffamily,
    >={Stealth[length=1.6mm]},
    bx/.style={draw, minimum height=6mm, align=center, inner sep=1pt},
    stp/.style={circle, draw, fill=white, inner sep=0.6pt, font=\scriptsize\sffamily\bfseries}]
  % GPU
  \draw[fill=ln-box] (0,0) rectangle (56,40);
  \node[anchor=north west, font=\scriptsize\sffamily\bfseries] at (0.5,39.5) {GPU};
  \foreach \k in {0,1,2} {
    \node[bx, fill=ln-accent!20, minimum width=14mm] at (8+\k*17,31) {SM\,\k\ +\,TLB};
  }
  \node[bx, fill=ln-caution!15, minimum width=50mm, minimum height=13mm] (gmmu) at (28,17)
    {};
  \node[anchor=north, inner sep=1pt] at (gmmu.north) {GMMU};
  \node[bx, fill=white, minimum width=14mm] at (11,15) {L2 TLB};
  \node[bx, fill=white, minimum width=14mm] at (28,15) {\iflangja{ウォーク}{walk}};
  \node[bx, fill=white, minimum width=14mm] at (45,15) {\iflangja{フォールト\\MSHR}{fault\\MSHRs}};
  \node[bx, fill=ln-tip!25, minimum width=50mm] (dmem) at (28,4.5) {\iflangja{デバイスメモリ（ページテーブル）}{device memory (page table)}};
  \foreach \k in {0,1,2} { \draw[->] (8+\k*17,28) -- (8+\k*17,23.5); }
  % PCIe
  \draw[fill=ln-muted!15] (60,14) rectangle (72,26);
  \node[rotate=90] at (66,20) {PCIe};
  % CPU
  \draw[fill=ln-box] (76,0) rectangle (116,40);
  \node[anchor=north west, font=\scriptsize\sffamily\bfseries] at (76.5,39.5) {CPU};
  \node[bx, fill=ln-accent!20, minimum width=34mm] (drv) at (96,26) {\iflangja{OS ＋ GPU ドライバ／\\ランタイム}{OS + GPU driver\\and runtime}};
  \node[bx, fill=ln-tip!25, minimum width=34mm] (hmem) at (96,7) {\iflangja{ホストメモリ}{host memory}};
  \draw[<->] (drv) -- (hmem);
  % fault and reply
  \draw[->, thick, ln-caution] (53,24) -- (80,29);
  \node[stp] at (49,25) {1};
  \node[stp] at (66,30) {2};
  \node[anchor=south, text=ln-caution, align=center] at (66,31.8) {\iflangja{ページフォールト}{page fault}};
  \node[stp] at (106,17) {3};
  \draw[->, thick, ln-tip] (80,22) -- (53,11);
  \node[stp] at (66,10) {4};
  \node[anchor=north, text=ln-tip, align=center] at (66,8.5) {\iflangja{新しいマッピング\\とデータ}{new mapping\\and data}};
\end{tikzpicture}

  \caption{A far fault on a discrete GPU: (1) the GMMU finds no valid
    PTE and records the fault; (2) the fault crosses PCIe; (3) the OS
    and the GPU driver allocate and migrate the page; (4) the mapping
    returns and the access is replayed.}
  \label{fig:l07-fault-path}
\end{figure}

Both processors reach memory through a translation unit for devices.

\begin{definition}[IOMMU and DMA]
  \label{def:l07-iommu}
  \term{DMA (direct memory access)}[DMA] is the transfer of data
  between a device and memory without the CPU executing the copy.  An
  \term{IOMMU}[IOMMU] (input--output memory management unit) translates
  the addresses that devices use for DMA, so that devices can work
  with virtual instead of physical addresses (I/O
  virtualisation).\index{direct memory access|see{DMA (direct memory access)}}
\end{definition}

\margin{Ganguly et al.\ give a far-fault handling latency of typically
\SI{45}{\micro\second} on Pascal GPUs, before the PCIe transfer of the
data \cite{ganguly2019interplay}.}%
GPU memory was once managed only through the IOMMU, as the memory of a
peripheral device.  With a shared address space the CPU takes part in
every mapping change, and this is the cost of on-demand GPU memory:
page management is performed by the host, and fetching a mapping from
the CPU costs a PCIe round trip plus the CPU's interrupt service
routine, often several tens of microseconds.\sidenote{The driver does not handle one fault at a time: it collects
faults in a buffer and processes them in batches, and faults raised
while a batch is processed wait for the next one.  Kim et al.\
\cite{kim2020batch} show that this serialisation hurts throughput and
propose larger batches (by switching thread-block contexts so that more
faults are raised together) and eviction overlapped with migration: 2$\times$ average speed-up over
state-of-the-art prefetching, batches 2.27$\times$ larger and 51\,\%
fewer, batch processing time 27\,\% shorter.}  Three remedies follow.
\begin{itemize}
  \item \emph{Prefetch pages}: predict the pages that will be needed
    and migrate them before they fault.
  \item \emph{Move page-mapping management to the GPU side}, so that
    a fault does not need the host.
  \item \emph{Let CPU and GPU cooperate}: the CPU manages large chunks
    of memory, the GPU the small pages inside a chunk it has been
    given.
\end{itemize}

Prefetching is the remedy that NVIDIA's driver implements
\cite{ganguly2019interplay}.

\subsection{Page prefetching}

\begin{definition}[Prefetching]
  \label{def:l07-prefetching}
  Fetching data before it is requested, on the strength of a
  prediction of future addresses, is called
  \term{prefetching}[プリフェッチ].
  Applied to the pages of a unified address space it is
  \term{page prefetching}[ページプリフェッチ]: migrating pages to device
  memory before the GPU faults on them.
\end{definition}

A sequential predictor that sees the addresses \texttt{0x40},
\texttt{0x80}, \texttt{0xC0} predicts \texttt{0x100} and
\texttt{0x140} next.  A vector addition that reads \texttt{a[idx]} and
\texttt{b[idx]} with consecutive indices produces exactly such a
sequence, and page addresses of such an array are as predictable: a
sequential page predictor fetches page $p+1$ when page $p$ faults.

\begin{marginfigure}
  % Page prefetching under out-of-order thread blocks (Lesson 7): an
% array spans nine pages; block 0 reads pages 0--2, block 1 pages
% 3--5, block 2 pages 6--8.  The blocks run in the order 2, 0, 1, so
% the faults arrive as 6 7 8 0 1 2 3 4 5 and a predictor of "next
% page" mispredicts at every block boundary.
% Redraws the outline figure m7l2_preteching (page numbers made
% consistent).
% Usage: % Page prefetching under out-of-order thread blocks (Lesson 7): an
% array spans nine pages; block 0 reads pages 0--2, block 1 pages
% 3--5, block 2 pages 6--8.  The blocks run in the order 2, 0, 1, so
% the faults arrive as 6 7 8 0 1 2 3 4 5 and a predictor of "next
% page" mispredicts at every block boundary.
% Redraws the outline figure m7l2_preteching (page numbers made
% consistent).
% Usage: % Page prefetching under out-of-order thread blocks (Lesson 7): an
% array spans nine pages; block 0 reads pages 0--2, block 1 pages
% 3--5, block 2 pages 6--8.  The blocks run in the order 2, 0, 1, so
% the faults arrive as 6 7 8 0 1 2 3 4 5 and a predictor of "next
% page" mispredicts at every block boundary.
% Redraws the outline figure m7l2_preteching (page numbers made
% consistent).
% Usage: \input{figures/l07-block-pages}  (in a marginfigure)
% Width: margin (50 mm); the drawing is about 48 mm wide.
\begin{tikzpicture}[x=1mm, y=1mm, font=\scriptsize\sffamily]
  \foreach \p in {0,...,8} {
    \pgfmathtruncatemacro{\b}{\p/3}
    \ifcase\b \def\col{ln-accent!25}\or \def\col{ln-caution!20}\or \def\col{ln-tip!25}\fi
    \draw[fill=\col] (\p*5.2,10) rectangle ++(5.2,6) node[midway] {\p};
  }
  \node[anchor=south west, text=ln-muted] at (0,16) {\iflangja{ページ}{page}};
  \foreach \b/\col in {0/ln-accent!25, 1/ln-caution!20, 2/ln-tip!25} {
    \draw[fill=\col] (\b*15.6+0.5,2) rectangle ++(14.6,5)
      node[midway] {\iflangja{ブロック}{block}~\b};
  }
  \node[anchor=north west, text=ln-muted, align=left] at (0,1)
    {\iflangja{実行順：2, 0, 1}{run order: 2, 0, 1}};
\end{tikzpicture}
  (in a marginfigure)
% Width: margin (50 mm); the drawing is about 48 mm wide.
\begin{tikzpicture}[x=1mm, y=1mm, font=\scriptsize\sffamily]
  \foreach \p in {0,...,8} {
    \pgfmathtruncatemacro{\b}{\p/3}
    \ifcase\b \def\col{ln-accent!25}\or \def\col{ln-caution!20}\or \def\col{ln-tip!25}\fi
    \draw[fill=\col] (\p*5.2,10) rectangle ++(5.2,6) node[midway] {\p};
  }
  \node[anchor=south west, text=ln-muted] at (0,16) {\iflangja{ページ}{page}};
  \foreach \b/\col in {0/ln-accent!25, 1/ln-caution!20, 2/ln-tip!25} {
    \draw[fill=\col] (\b*15.6+0.5,2) rectangle ++(14.6,5)
      node[midway] {\iflangja{ブロック}{block}~\b};
  }
  \node[anchor=north west, text=ln-muted, align=left] at (0,1)
    {\iflangja{実行順：2, 0, 1}{run order: 2, 0, 1}};
\end{tikzpicture}
  (in a marginfigure)
% Width: margin (50 mm); the drawing is about 48 mm wide.
\begin{tikzpicture}[x=1mm, y=1mm, font=\scriptsize\sffamily]
  \foreach \p in {0,...,8} {
    \pgfmathtruncatemacro{\b}{\p/3}
    \ifcase\b \def\col{ln-accent!25}\or \def\col{ln-caution!20}\or \def\col{ln-tip!25}\fi
    \draw[fill=\col] (\p*5.2,10) rectangle ++(5.2,6) node[midway] {\p};
  }
  \node[anchor=south west, text=ln-muted] at (0,16) {\iflangja{ページ}{page}};
  \foreach \b/\col in {0/ln-accent!25, 1/ln-caution!20, 2/ln-tip!25} {
    \draw[fill=\col] (\b*15.6+0.5,2) rectangle ++(14.6,5)
      node[midway] {\iflangja{ブロック}{block}~\b};
  }
  \node[anchor=north west, text=ln-muted, align=left] at (0,1)
    {\iflangja{実行順：2, 0, 1}{run order: 2, 0, 1}};
\end{tikzpicture}

  \caption{Nine pages read by three thread blocks that run in the
    order 2, 0, 1: the faults arrive as 6, 7, 8, 0, 1, \ldots}
  \label{fig:l07-block-pages}
\end{marginfigure}

On a GPU the prediction breaks down, because thread blocks do not run
in order (Lesson~3).  In \cref{fig:l07-block-pages} each block reads
three consecutive pages, but the blocks run in the order 2, 0, 1, and
the faults arrive as 6, 7, 8, 0, 1, 2, 3, 4, 5: after page 8 a
\enquote{next page} predictor fetches page 9, which no block reads, and
the jump back to page 0 is a fault it could not foresee.  Spatial locality is there, in the neighbourhood of a
fault, but not in the order of the faults.

\begin{definition}[Tree-based neighborhood prefetching]
  \label{def:l07-tbnp}
  The policy of NVIDIA's unified-memory driver, called
  \term{tree-based neighborhood prefetching}[ツリーベース近傍プリフェッチ]
  (TBNp), divides each
  \SI{2}{\mebi\byte} region of a managed allocation into two
  \SI{1}{\mebi\byte} halves and those into \SI{64}{\kibi\byte} nodes.
  A \SI{4}{\kibi\byte} fault migrates the whole \SI{64}{\kibi\byte}
  node, 16 pages; when several nodes under one parent have faulted, the
  parent node, up to \SI{1}{\mebi\byte}, is fetched.\index{TBNp|see{tree-based neighborhood prefetching}}
\end{definition}

The tree captures spatial locality around each fault, whatever the
order of the faults, and bounds the size of a single prefetch.\caution{Lecture vs paper: Ganguly et al.\ \cite{ganguly2019interplay}
measured the exact rule on a GTX 1080 Ti.  The tree is a full binary
tree with \SI{64}{\kibi\byte} leaves, so it has the levels 128, 256 and
\SI{512}{\kibi\byte} in between; a node is filled up when its
\enquote{to-be valid size} is \enquote{strictly greater than 50\,\%}
of its capacity, so a prefetch ranges from \SI{64}{\kibi\byte} to
\SI{1}{\mebi\byte}.}

\begin{example}
  \label{ex:l07-tbnp}
  Apply the rule of Ganguly et al.\ to far faults in the leaves 0, 1,
  2 and 4 of a \SI{2}{\mebi\byte} region, in this order
  (\cref{fig:l07-tbnp-tree}).  After each fault the GMMU adds up the
  valid size of every ancestor; a node more than half valid has its
  emptier child filled up to the size of the fuller one.
  \begin{center}
    \small
    \begin{tabular}{@{}c>{\raggedright\arraybackslash}p{58mm}c@{}}
      \toprule
      fault & ancestors above 50\,\% & migrated \\
      \midrule
      leaf 0 & none (128-KiB node: 64 of 128) & \SI{64}{\kibi\byte} \\
      leaf 1 & 128-KiB node full, nothing to fill (256-KiB node: 128 of 256) & \SI{64}{\kibi\byte} \\
      leaf 2 & 256-KiB node: 192 of 256; fill leaf 3 & \SI{128}{\kibi\byte} \\
      leaf 4 & 512-KiB node: 320 of 512; fill leaves 5--7 & \SI{256}{\kibi\byte} \\
      \bottomrule
    \end{tabular}
  \end{center}
  Four faults bring in \SI{512}{\kibi\byte}, eight leaves, instead of
  the \SI{16}{\kibi\byte} that four on-demand \SI{4}{\kibi\byte}
  migrations would.  Leaf 0 alone did not trigger a prefetch beyond
  its own node: at exactly 50\,\% a node is not filled.
\end{example}

\begin{figure}[tb]
  \centering
  % Tree-based neighbourhood prefetching (Lesson 7): a 2 MiB region of a
% managed allocation as a full binary tree whose 32 leaves are 64 KiB
% basic blocks (16 pages of 4 KiB).  Shaded leaves show the end state
% of the worked example in the text: far faults in leaves 0, 1, 2 and
% 4 (red) and the leaves prefetched by the >50 % rule (green).
% Redraws the outline figure m7l2_treebase, with the intermediate
% levels of Ganguly et al. (2019) drawn in.
% Usage: % Tree-based neighbourhood prefetching (Lesson 7): a 2 MiB region of a
% managed allocation as a full binary tree whose 32 leaves are 64 KiB
% basic blocks (16 pages of 4 KiB).  Shaded leaves show the end state
% of the worked example in the text: far faults in leaves 0, 1, 2 and
% 4 (red) and the leaves prefetched by the >50 % rule (green).
% Redraws the outline figure m7l2_treebase, with the intermediate
% levels of Ganguly et al. (2019) drawn in.
% Usage: % Tree-based neighbourhood prefetching (Lesson 7): a 2 MiB region of a
% managed allocation as a full binary tree whose 32 leaves are 64 KiB
% basic blocks (16 pages of 4 KiB).  Shaded leaves show the end state
% of the worked example in the text: far faults in leaves 0, 1, 2 and
% 4 (red) and the leaves prefetched by the >50 % rule (green).
% Redraws the outline figure m7l2_treebase, with the intermediate
% levels of Ganguly et al. (2019) drawn in.
% Usage: \input{figures/l07-tbnp-tree}
% Width: body (120 mm); the drawing is about 117 mm wide.
\begin{tikzpicture}[x=1mm, y=1mm, font=\scriptsize\sffamily]
  \def\w{3.3}
  % level labels
  \foreach \h/\lab in {0/64\,KiB, 1/128\,KiB, 2/256\,KiB, 3/512\,KiB, 4/1\,MiB, 5/2\,MiB} {
    \node[anchor=east, text=ln-muted] at (11,\h*6+2) {\lab};
  }
  % internal nodes and edges
  \foreach \h in {1,...,5} {
    \pgfmathtruncatemacro{\n}{2^(5-\h)-1}
    \pgfmathsetmacro{\span}{2^\h*\w}
    \foreach \j in {0,...,\n} {
      \pgfmathsetmacro{\x}{12+\j*\span+\span/2}
      \pgfmathsetmacro{\xl}{12+\j*\span+\span/4}
      \pgfmathsetmacro{\xr}{12+\j*\span+3*\span/4}
      \ifnum\h=1
        \draw[ln-rule] (\x,\h*6+2) -- (\xl,4.5);
        \draw[ln-rule] (\x,\h*6+2) -- (\xr,4.5);
      \else
        \draw[ln-rule] (\x,\h*6+2) -- (\xl,\h*6-4);
        \draw[ln-rule] (\x,\h*6+2) -- (\xr,\h*6-4);
      \fi
      \fill[ln-accent] (\x,\h*6+2) circle (0.7);
    }
  }
  % leaves
  \foreach \i in {0,...,31} {
    \def\col{white}
    \ifnum\i<3 \def\col{ln-caution!35}\fi
    \ifnum\i=4 \def\col{ln-caution!35}\fi
    \ifnum\i=3 \def\col{ln-tip!40}\fi
    \ifnum\i>4 \ifnum\i<8 \def\col{ln-tip!40}\fi\fi
    \draw[fill=\col] (12+\i*\w,-0.5) rectangle ++(\w,5);
  }
  \foreach \i in {0,4,8,12,16,20,24,28} {
    \node[anchor=north, text=ln-muted] at (12+\i*\w+\w/2,-0.6) {\i};
  }
  % legend
  \draw[fill=ln-caution!35] (80,30) rectangle ++(3,3);
  \node[anchor=west] at (83.5,31.5) {\iflangja{ファーフォールト}{far fault}};
  \draw[fill=ln-tip!40] (98,30) rectangle ++(3,3);
  \node[anchor=west] at (101.5,31.5) {\iflangja{プリフェッチ}{prefetched}};
\end{tikzpicture}

% Width: body (120 mm); the drawing is about 117 mm wide.
\begin{tikzpicture}[x=1mm, y=1mm, font=\scriptsize\sffamily]
  \def\w{3.3}
  % level labels
  \foreach \h/\lab in {0/64\,KiB, 1/128\,KiB, 2/256\,KiB, 3/512\,KiB, 4/1\,MiB, 5/2\,MiB} {
    \node[anchor=east, text=ln-muted] at (11,\h*6+2) {\lab};
  }
  % internal nodes and edges
  \foreach \h in {1,...,5} {
    \pgfmathtruncatemacro{\n}{2^(5-\h)-1}
    \pgfmathsetmacro{\span}{2^\h*\w}
    \foreach \j in {0,...,\n} {
      \pgfmathsetmacro{\x}{12+\j*\span+\span/2}
      \pgfmathsetmacro{\xl}{12+\j*\span+\span/4}
      \pgfmathsetmacro{\xr}{12+\j*\span+3*\span/4}
      \ifnum\h=1
        \draw[ln-rule] (\x,\h*6+2) -- (\xl,4.5);
        \draw[ln-rule] (\x,\h*6+2) -- (\xr,4.5);
      \else
        \draw[ln-rule] (\x,\h*6+2) -- (\xl,\h*6-4);
        \draw[ln-rule] (\x,\h*6+2) -- (\xr,\h*6-4);
      \fi
      \fill[ln-accent] (\x,\h*6+2) circle (0.7);
    }
  }
  % leaves
  \foreach \i in {0,...,31} {
    \def\col{white}
    \ifnum\i<3 \def\col{ln-caution!35}\fi
    \ifnum\i=4 \def\col{ln-caution!35}\fi
    \ifnum\i=3 \def\col{ln-tip!40}\fi
    \ifnum\i>4 \ifnum\i<8 \def\col{ln-tip!40}\fi\fi
    \draw[fill=\col] (12+\i*\w,-0.5) rectangle ++(\w,5);
  }
  \foreach \i in {0,4,8,12,16,20,24,28} {
    \node[anchor=north, text=ln-muted] at (12+\i*\w+\w/2,-0.6) {\i};
  }
  % legend
  \draw[fill=ln-caution!35] (80,30) rectangle ++(3,3);
  \node[anchor=west] at (83.5,31.5) {\iflangja{ファーフォールト}{far fault}};
  \draw[fill=ln-tip!40] (98,30) rectangle ++(3,3);
  \node[anchor=west] at (101.5,31.5) {\iflangja{プリフェッチ}{prefetched}};
\end{tikzpicture}

% Width: body (120 mm); the drawing is about 117 mm wide.
\begin{tikzpicture}[x=1mm, y=1mm, font=\scriptsize\sffamily]
  \def\w{3.3}
  % level labels
  \foreach \h/\lab in {0/64\,KiB, 1/128\,KiB, 2/256\,KiB, 3/512\,KiB, 4/1\,MiB, 5/2\,MiB} {
    \node[anchor=east, text=ln-muted] at (11,\h*6+2) {\lab};
  }
  % internal nodes and edges
  \foreach \h in {1,...,5} {
    \pgfmathtruncatemacro{\n}{2^(5-\h)-1}
    \pgfmathsetmacro{\span}{2^\h*\w}
    \foreach \j in {0,...,\n} {
      \pgfmathsetmacro{\x}{12+\j*\span+\span/2}
      \pgfmathsetmacro{\xl}{12+\j*\span+\span/4}
      \pgfmathsetmacro{\xr}{12+\j*\span+3*\span/4}
      \ifnum\h=1
        \draw[ln-rule] (\x,\h*6+2) -- (\xl,4.5);
        \draw[ln-rule] (\x,\h*6+2) -- (\xr,4.5);
      \else
        \draw[ln-rule] (\x,\h*6+2) -- (\xl,\h*6-4);
        \draw[ln-rule] (\x,\h*6+2) -- (\xr,\h*6-4);
      \fi
      \fill[ln-accent] (\x,\h*6+2) circle (0.7);
    }
  }
  % leaves
  \foreach \i in {0,...,31} {
    \def\col{white}
    \ifnum\i<3 \def\col{ln-caution!35}\fi
    \ifnum\i=4 \def\col{ln-caution!35}\fi
    \ifnum\i=3 \def\col{ln-tip!40}\fi
    \ifnum\i>4 \ifnum\i<8 \def\col{ln-tip!40}\fi\fi
    \draw[fill=\col] (12+\i*\w,-0.5) rectangle ++(\w,5);
  }
  \foreach \i in {0,4,8,12,16,20,24,28} {
    \node[anchor=north, text=ln-muted] at (12+\i*\w+\w/2,-0.6) {\i};
  }
  % legend
  \draw[fill=ln-caution!35] (80,30) rectangle ++(3,3);
  \node[anchor=west] at (83.5,31.5) {\iflangja{ファーフォールト}{far fault}};
  \draw[fill=ln-tip!40] (98,30) rectangle ++(3,3);
  \node[anchor=west] at (101.5,31.5) {\iflangja{プリフェッチ}{prefetched}};
\end{tikzpicture}

  \caption{The prefetch tree of a \SI{2}{\mebi\byte} region: 32 leaves
    of \SI{64}{\kibi\byte} (16 pages each) under binary levels up to
    the root.  Shaded: the end state of \cref{ex:l07-tbnp}, with far
    faults in leaves 0, 1, 2 and 4 and the leaves prefetched.}
  \label{fig:l07-tbnp-tree}
\end{figure}

\subsection{Oversubscription and eviction}

\begin{definition}[Oversubscription and eviction]
  \label{def:l07-oversubscription}
  The allocation of more managed memory than the GPU's physical memory
  holds is called
  \term{memory oversubscription}[メモリのオーバーサブスクリプション].
  Pages must then leave device memory to make room, which is
  \term{page eviction}[ページ追い出し]; the policy decides which pages
  go.
\end{definition}

\margin{CPG §19.1.2.7 GPU Memory Oversubscription: unified memory
lets applications \enquote{allocate and share arrays larger than the
memory capacity of any individual processor in the system}
\cite{nvidia2024cudaguide}.}%
Prefetching and eviction interact.  Aggressive prefetching evicts
pages that are still needed, and an eviction policy that ignores the
prefetcher evicts scattered \SI{4}{\kibi\byte} pages, which breaks the
contiguous regions the prefetcher relies on.  The problem grows with
the memory demand of kernels.  One answer is
\term{pre-eviction}[事前追い出し]: together with the page that the
LRU list selects for eviction, the contiguous pages around it are
evicted in multiples of the \SI{64}{\kibi\byte} blocks in which the
prefetcher brought them in, so that free space appears in units the
prefetcher can use again.\margin{Ganguly et al.: a
tree-based pre-eviction combined with the hardware prefetcher gives
93\,\% and 18.5\,\% average speed-ups over LRU replacement of
\SI{4}{\kibi\byte} and \SI{2}{\mebi\byte} pages
\cite{ganguly2019interplay}.}

\begin{keypoint}
  A shared address space removes the explicit copies, not the data
  movement: on a discrete GPU every first touch of a page is a far
  fault that crosses PCIe to the host and back, tens of microseconds.
  Prefetching around each fault (the tree-based scheme) amortises the
  round trips despite the out-of-order execution of thread blocks;
  under oversubscription the eviction policy must respect the units
  the prefetcher works in.
\end{keypoint}

\section{Warp scheduling}
\label{sec:l07-scheduling}

The pipeline walk-through of Lesson~4 assumed one ready warp at a
time.\source{M7, GTO vs round-robin; A18 ch.~3; K12 §1.4.1 GPU
Pipeline; CPG §4.2 Hardware Multithreading; \cite{rogers2012ccws,
narasiman2011twolevel,gebhart2011energy}.}  Usually several warps have
all their operands ready in the scoreboard at once, and one of them
must be chosen.  This section defines the basic policies, shows how the
choice affects latency hiding and the L1 cache, and presents research
schedulers built on these observations.

\begin{definition}[Warp scheduler]
  \label{def:l07-warp-scheduler}
  The \term{warp scheduler}[ワープスケジューラ] is the pipeline stage
  that, every cycle, selects one warp among those whose next
  instruction is ready and issues that instruction.
\end{definition}

\margin{CPG §16.6.1 (Compute Capability 7.x, Architecture): an SM has
4 warp schedulers; it distributes its warps statically among them, and
each issues one instruction per cycle for one of its ready warps
\cite{nvidia2024cudaguide}.  With the 64 resident warps of CC 7.0
(CPG Table 21; 32 on CC 7.5) that is up to 16 warps per scheduler (our
arithmetic).}%
\subsection{From instruction scheduling to warp scheduling}

An out-of-order CPU has a similar unit.  Its instruction scheduler
selects among the ready instructions of a scheduling queue.  The
default policy is \term{oldest-first}[最古優先]: among the ready
instructions, the one that entered the queue first issues.
Alternatives give priority to \emph{critical} instructions: loads
first, so that memory requests start as soon as their addresses are
known, or instructions with many dependants first, so that more
instructions become ready sooner.

The GPU warp scheduler is in order within each warp: a warp's
instructions issue in program order.  Across warps it may pick any
ready warp, so the instruction stream of the SM as a whole looks out of
order.  Four warps that each run the instructions I1, I2, I3 may, for
example, issue as
\begin{center}
  \small
  I1\,w1, I1\,w2, I2\,w1, I1\,w3, I2\,w2, I3\,w1, I1\,w4, I2\,w3,
  I2\,w4, I3\,w2, I3\,w4, I3\,w3,
\end{center}
which keeps the order I1, I2, I3 inside every warp.\margin{K12 §1.4.1:
the scheduler is in order but \enquote{can select any warps that are
ready}, so the SM's execution \enquote{might look like an out-of-order
execution} \cite{kim2012gpgpu}.}  Which ready warp goes first is the
policy.

\subsection{Round-robin and greedy-then-oldest}

\begin{definition}[Round-robin and GTO]
  \label{def:l07-rr-gto}
  A \term{round-robin (RR)}[ラウンドロビン] scheduler gives the warps
  turns in a fixed cyclic order, one instruction each.  In the
  \term{loose round-robin (LRR)}[ルーズラウンドロビン] form, a warp
  that cannot issue in its turn is skipped and the next one in order
  gets the chance.  A
  \term{greedy-then-oldest (GTO)}[GTOスケジューリング] scheduler keeps
  issuing from one warp until it stalls on a long-latency event (a
  cache miss, a branch, a long tensor-core operation), then switches to
  the oldest ready warp, the age of a warp being the time it was
  assigned to the SM.\index{RR|see{round-robin (RR)}}\index{LRR|see{loose round-robin (LRR)}}\index{GTO|see{greedy-then-oldest (GTO)}}
\end{definition}

\margin{Rogers et al.\ define both baselines: GTO \enquote{runs a
single wavefront until it stalls then picks the oldest ready
wavefront}, the age being \enquote{the time it is assigned to the
core} \cite{rogers2012ccws}.}%
The two policies differ in when the warps reach their memory
instructions.  Under round-robin all warps advance together and reach
their loads at nearly the same time; under GTO one warp runs ahead and
starts its load early, and the others follow one by one.

\begin{example}
  \label{ex:l07-rr-gto}
  Four warps run the same program: two arithmetic instructions
  \texttt{c1}, \texttt{c2}, a load \texttt{m3} that misses in the
  cache and returns after 10 cycles, and three arithmetic instructions
  \texttt{c4}--\texttt{c6}, of which \texttt{c4} uses the loaded
  value.\margin{Notation: one instruction issues per cycle; a miss
issued in cycle $t$ lets the dependent instruction issue in cycle
$t+10$.  Our numbers.}  \Cref{fig:l07-rr-gto} shows the issue slots.
  \begin{itemize}
    \item \emph{Round-robin}: cycles 0--7 issue \texttt{c1} and
      \texttt{c2} of all warps, cycles 8--11 the four loads.  Every
      warp now waits for its load; the first returns in cycle 18, so
      cycles 12--17 are stalls.  The 12 remaining instructions take
      cycles 18--29: 30 cycles, 6 of them stalls.
    \item \emph{GTO}: warp 1 issues \texttt{c1}, \texttt{c2},
      \texttt{m3} in cycles 0--2 and stalls; warp 2 follows in cycles
      3--5, warp 3 in 6--8, warp 4 in 9--11.  Warp 1's load returns
      in cycle 12, one cycle after warp 4's load issues, and each warp then
      finishes its three instructions as its data arrives: 24 cycles,
      no stall.
  \end{itemize}
  The loads under GTO are spread out, and the other warps' arithmetic
  hides each one.
\end{example}

\begin{figure}[tb]
  \centering
  % Round-robin versus greedy-then-oldest issue (Lesson 7), the worked
% example of the text: four warps each run c1 c2 m3 c4 c5 c6, where m3
% is a load that misses (latency 10 cycles) and c4 needs its value.
% One instruction issues per cycle; the number in a cell is the warp,
% M marks the load, grey cells are stall cycles, and the bars below a
% row show each load's latency.
% Redraws the outline figure m7l3_rr_gto with the text's own program.
% Usage: % Round-robin versus greedy-then-oldest issue (Lesson 7), the worked
% example of the text: four warps each run c1 c2 m3 c4 c5 c6, where m3
% is a load that misses (latency 10 cycles) and c4 needs its value.
% One instruction issues per cycle; the number in a cell is the warp,
% M marks the load, grey cells are stall cycles, and the bars below a
% row show each load's latency.
% Redraws the outline figure m7l3_rr_gto with the text's own program.
% Usage: % Round-robin versus greedy-then-oldest issue (Lesson 7), the worked
% example of the text: four warps each run c1 c2 m3 c4 c5 c6, where m3
% is a load that misses (latency 10 cycles) and c4 needs its value.
% One instruction issues per cycle; the number in a cell is the warp,
% M marks the load, grey cells are stall cycles, and the bars below a
% row show each load's latency.
% Redraws the outline figure m7l3_rr_gto with the text's own program.
% Usage: \input{figures/l07-rr-gto}
% Width: body (120 mm); the drawing is about 118 mm wide.
\begin{tikzpicture}[x=3.5mm, y=1mm, font=\scriptsize\sffamily]
  \colorlet{l07d1}{ln-accent}\colorlet{l07d2}{ln-caution}%
  \colorlet{l07d3}{ln-tip}\colorlet{l07d4}{ln-source}%
  % \lsevenissue{y}{list of w/m entries; 0 = stall}
  \newcommand{\lsevenissue}[2]{%
    \foreach \w/\m [count=\c from 0] in {#2} {
      \ifnum\w=0
        \draw[fill=ln-rule!40, draw=ln-rule] (\c+3,#1) rectangle ++(1,5);
      \else
        \draw[fill=l07d\w!22] (\c+3,#1) rectangle ++(1,5)
          node[midway] {\ifnum\m=1 M\else\w\fi};
      \fi
    }}
  % \lsevenlat{y}{issue cycle}{warp}
  \newcommand{\lsevenlat}[3]{%
    \draw[|-|, thick, draw=l07d#3] (#2+3.5,#1) -- (#2+13.5,#1);}
  % round-robin
  \iflangja{\node[anchor=east, align=right, inner sep=0.3pt] at (2.8,30.5) {ラウンド\\ロビン};}%
    {\node[anchor=east, align=right] at (2.8,30.5) {round-\\robin};}
  \lsevenissue{28}{1/0,2/0,3/0,4/0,1/0,2/0,3/0,4/0,1/1,2/1,3/1,4/1,0/0,0/0,0/0,0/0,0/0,0/0,1/0,2/0,3/0,4/0,1/0,2/0,3/0,4/0,1/0,2/0,3/0,4/0}
  \lsevenlat{26}{8}{1}\lsevenlat{24.5}{9}{2}\lsevenlat{23}{10}{3}\lsevenlat{21.5}{11}{4}
  % GTO
  \node[anchor=east, align=right] at (2.8,11.5) {GTO};
  \lsevenissue{9}{1/0,1/0,1/1,2/0,2/0,2/1,3/0,3/0,3/1,4/0,4/0,4/1,1/0,1/0,1/0,2/0,2/0,2/0,3/0,3/0,3/0,4/0,4/0,4/0}
  \lsevenlat{7}{2}{1}\lsevenlat{5.5}{5}{2}\lsevenlat{4}{8}{3}\lsevenlat{2.5}{11}{4}
  % cycle axis
  \foreach \c in {0,5,10,15,20,25,30} {
    \draw[ln-muted] (\c+3,0) -- ++(0,-1) node[below, inner sep=1pt] {\c};
  }
  \draw[->, ln-muted] (3,0) -- (34,0) node[above left, inner sep=1pt] {\iflangja{サイクル}{cycle}};
\end{tikzpicture}

% Width: body (120 mm); the drawing is about 118 mm wide.
\begin{tikzpicture}[x=3.5mm, y=1mm, font=\scriptsize\sffamily]
  \colorlet{l07d1}{ln-accent}\colorlet{l07d2}{ln-caution}%
  \colorlet{l07d3}{ln-tip}\colorlet{l07d4}{ln-source}%
  % \lsevenissue{y}{list of w/m entries; 0 = stall}
  \newcommand{\lsevenissue}[2]{%
    \foreach \w/\m [count=\c from 0] in {#2} {
      \ifnum\w=0
        \draw[fill=ln-rule!40, draw=ln-rule] (\c+3,#1) rectangle ++(1,5);
      \else
        \draw[fill=l07d\w!22] (\c+3,#1) rectangle ++(1,5)
          node[midway] {\ifnum\m=1 M\else\w\fi};
      \fi
    }}
  % \lsevenlat{y}{issue cycle}{warp}
  \newcommand{\lsevenlat}[3]{%
    \draw[|-|, thick, draw=l07d#3] (#2+3.5,#1) -- (#2+13.5,#1);}
  % round-robin
  \iflangja{\node[anchor=east, align=right, inner sep=0.3pt] at (2.8,30.5) {ラウンド\\ロビン};}%
    {\node[anchor=east, align=right] at (2.8,30.5) {round-\\robin};}
  \lsevenissue{28}{1/0,2/0,3/0,4/0,1/0,2/0,3/0,4/0,1/1,2/1,3/1,4/1,0/0,0/0,0/0,0/0,0/0,0/0,1/0,2/0,3/0,4/0,1/0,2/0,3/0,4/0,1/0,2/0,3/0,4/0}
  \lsevenlat{26}{8}{1}\lsevenlat{24.5}{9}{2}\lsevenlat{23}{10}{3}\lsevenlat{21.5}{11}{4}
  % GTO
  \node[anchor=east, align=right] at (2.8,11.5) {GTO};
  \lsevenissue{9}{1/0,1/0,1/1,2/0,2/0,2/1,3/0,3/0,3/1,4/0,4/0,4/1,1/0,1/0,1/0,2/0,2/0,2/0,3/0,3/0,3/0,4/0,4/0,4/0}
  \lsevenlat{7}{2}{1}\lsevenlat{5.5}{5}{2}\lsevenlat{4}{8}{3}\lsevenlat{2.5}{11}{4}
  % cycle axis
  \foreach \c in {0,5,10,15,20,25,30} {
    \draw[ln-muted] (\c+3,0) -- ++(0,-1) node[below, inner sep=1pt] {\c};
  }
  \draw[->, ln-muted] (3,0) -- (34,0) node[above left, inner sep=1pt] {\iflangja{サイクル}{cycle}};
\end{tikzpicture}

% Width: body (120 mm); the drawing is about 118 mm wide.
\begin{tikzpicture}[x=3.5mm, y=1mm, font=\scriptsize\sffamily]
  \colorlet{l07d1}{ln-accent}\colorlet{l07d2}{ln-caution}%
  \colorlet{l07d3}{ln-tip}\colorlet{l07d4}{ln-source}%
  % \lsevenissue{y}{list of w/m entries; 0 = stall}
  \newcommand{\lsevenissue}[2]{%
    \foreach \w/\m [count=\c from 0] in {#2} {
      \ifnum\w=0
        \draw[fill=ln-rule!40, draw=ln-rule] (\c+3,#1) rectangle ++(1,5);
      \else
        \draw[fill=l07d\w!22] (\c+3,#1) rectangle ++(1,5)
          node[midway] {\ifnum\m=1 M\else\w\fi};
      \fi
    }}
  % \lsevenlat{y}{issue cycle}{warp}
  \newcommand{\lsevenlat}[3]{%
    \draw[|-|, thick, draw=l07d#3] (#2+3.5,#1) -- (#2+13.5,#1);}
  % round-robin
  \iflangja{\node[anchor=east, align=right, inner sep=0.3pt] at (2.8,30.5) {ラウンド\\ロビン};}%
    {\node[anchor=east, align=right] at (2.8,30.5) {round-\\robin};}
  \lsevenissue{28}{1/0,2/0,3/0,4/0,1/0,2/0,3/0,4/0,1/1,2/1,3/1,4/1,0/0,0/0,0/0,0/0,0/0,0/0,1/0,2/0,3/0,4/0,1/0,2/0,3/0,4/0,1/0,2/0,3/0,4/0}
  \lsevenlat{26}{8}{1}\lsevenlat{24.5}{9}{2}\lsevenlat{23}{10}{3}\lsevenlat{21.5}{11}{4}
  % GTO
  \node[anchor=east, align=right] at (2.8,11.5) {GTO};
  \lsevenissue{9}{1/0,1/0,1/1,2/0,2/0,2/1,3/0,3/0,3/1,4/0,4/0,4/1,1/0,1/0,1/0,2/0,2/0,2/0,3/0,3/0,3/0,4/0,4/0,4/0}
  \lsevenlat{7}{2}{1}\lsevenlat{5.5}{5}{2}\lsevenlat{4}{8}{3}\lsevenlat{2.5}{11}{4}
  % cycle axis
  \foreach \c in {0,5,10,15,20,25,30} {
    \draw[ln-muted] (\c+3,0) -- ++(0,-1) node[below, inner sep=1pt] {\c};
  }
  \draw[->, ln-muted] (3,0) -- (34,0) node[above left, inner sep=1pt] {\iflangja{サイクル}{cycle}};
\end{tikzpicture}

  \caption{Issue slots of \cref{ex:l07-rr-gto}: the number is the warp,
    M its load, grey a stall cycle; the bars below each row are the
    10-cycle load latencies.  Round-robin issues the four loads back
    to back and stalls; GTO staggers them.}
  \label{fig:l07-rr-gto}
\end{figure}

GTO is not always better: the gain depends on the program.  In the
example it comes from the arithmetic after the load.  A program whose
warps all stall at the same barrier right after their loads gains
nothing; with a load latency longer than the arithmetic of all other
warps together GTO stalls too; and when warps share data, issuing them
together under round-robin may let one warp's miss serve the others.

\subsection{Two-level scheduling}

A scheduler that searches all ready warps every cycle compares tens of
candidates for one issue slot, and this search costs power.

\begin{marginfigure}
  % Two-level warp scheduling (Lesson 7): the scheduler searches only the
% small active set; a warp that meets a long-latency operation moves to
% the pending set, and a pending warp whose operands are ready refills
% the active set.
% Redraws the outline figure m7l3_two_lvl (right half).
% Usage: % Two-level warp scheduling (Lesson 7): the scheduler searches only the
% small active set; a warp that meets a long-latency operation moves to
% the pending set, and a pending warp whose operands are ready refills
% the active set.
% Redraws the outline figure m7l3_two_lvl (right half).
% Usage: % Two-level warp scheduling (Lesson 7): the scheduler searches only the
% small active set; a warp that meets a long-latency operation moves to
% the pending set, and a pending warp whose operands are ready refills
% the active set.
% Redraws the outline figure m7l3_two_lvl (right half).
% Usage: \input{figures/l07-two-level}  (in a marginfigure)
% Width: margin (50 mm); the drawing is about 46 mm wide.
\begin{tikzpicture}[x=1mm, y=1mm, font=\scriptsize\sffamily,
    >={Stealth[length=1.6mm]},
    wp/.style={draw, minimum width=8mm, minimum height=4.5mm, inner sep=1pt}]
  \node[anchor=south, text=ln-muted] at (8,38) {\iflangja{待機集合}{pending}};
  \foreach \w [count=\k from 0] in {4,5,6,7,8} {
    \node[wp, fill=ln-box] at (8,35-\k*5.5) {W\w};
  }
  \node[anchor=south, text=ln-muted] at (34,38) {\iflangja{アクティブ集合}{active}};
  \foreach \w [count=\k from 0] in {1,2,3} {
    \node[wp, fill=ln-accent!20] at (34,35-\k*5.5) {W\w};
  }
  \draw[->, ln-caution] (29.5,30) to[bend right=20] node[above, text=ln-caution, align=center] {\iflangja{長レイテンシ}{long}\\[-2pt]\iflangja{}{latency}} (12.5,30);
  \draw[->, ln-tip] (12.5,19) to[bend right=20] node[below, text=ln-tip] {\iflangja{準備完了}{ready}} (29.5,19);
  \node[draw, fill=ln-tip!25, minimum width=16mm, minimum height=5mm] (s) at (34,6) {\iflangja{スケジューラ}{scheduler}};
  \draw[->] (34,21.7) -- (s);
\end{tikzpicture}
  (in a marginfigure)
% Width: margin (50 mm); the drawing is about 46 mm wide.
\begin{tikzpicture}[x=1mm, y=1mm, font=\scriptsize\sffamily,
    >={Stealth[length=1.6mm]},
    wp/.style={draw, minimum width=8mm, minimum height=4.5mm, inner sep=1pt}]
  \node[anchor=south, text=ln-muted] at (8,38) {\iflangja{待機集合}{pending}};
  \foreach \w [count=\k from 0] in {4,5,6,7,8} {
    \node[wp, fill=ln-box] at (8,35-\k*5.5) {W\w};
  }
  \node[anchor=south, text=ln-muted] at (34,38) {\iflangja{アクティブ集合}{active}};
  \foreach \w [count=\k from 0] in {1,2,3} {
    \node[wp, fill=ln-accent!20] at (34,35-\k*5.5) {W\w};
  }
  \draw[->, ln-caution] (29.5,30) to[bend right=20] node[above, text=ln-caution, align=center] {\iflangja{長レイテンシ}{long}\\[-2pt]\iflangja{}{latency}} (12.5,30);
  \draw[->, ln-tip] (12.5,19) to[bend right=20] node[below, text=ln-tip] {\iflangja{準備完了}{ready}} (29.5,19);
  \node[draw, fill=ln-tip!25, minimum width=16mm, minimum height=5mm] (s) at (34,6) {\iflangja{スケジューラ}{scheduler}};
  \draw[->] (34,21.7) -- (s);
\end{tikzpicture}
  (in a marginfigure)
% Width: margin (50 mm); the drawing is about 46 mm wide.
\begin{tikzpicture}[x=1mm, y=1mm, font=\scriptsize\sffamily,
    >={Stealth[length=1.6mm]},
    wp/.style={draw, minimum width=8mm, minimum height=4.5mm, inner sep=1pt}]
  \node[anchor=south, text=ln-muted] at (8,38) {\iflangja{待機集合}{pending}};
  \foreach \w [count=\k from 0] in {4,5,6,7,8} {
    \node[wp, fill=ln-box] at (8,35-\k*5.5) {W\w};
  }
  \node[anchor=south, text=ln-muted] at (34,38) {\iflangja{アクティブ集合}{active}};
  \foreach \w [count=\k from 0] in {1,2,3} {
    \node[wp, fill=ln-accent!20] at (34,35-\k*5.5) {W\w};
  }
  \draw[->, ln-caution] (29.5,30) to[bend right=20] node[above, text=ln-caution, align=center] {\iflangja{長レイテンシ}{long}\\[-2pt]\iflangja{}{latency}} (12.5,30);
  \draw[->, ln-tip] (12.5,19) to[bend right=20] node[below, text=ln-tip] {\iflangja{準備完了}{ready}} (29.5,19);
  \node[draw, fill=ln-tip!25, minimum width=16mm, minimum height=5mm] (s) at (34,6) {\iflangja{スケジューラ}{scheduler}};
  \draw[->] (34,21.7) -- (s);
\end{tikzpicture}

  \caption{Two-level scheduling: the scheduler searches only the
    active set.}
  \label{fig:l07-two-level}
\end{marginfigure}

\begin{definition}[Two-level warp scheduling]
  \label{def:l07-two-level}
  In \term{two-level warp scheduling}[2レベルワープスケジューリング]
  the warps are divided into an
  \term{active warp set / pending warp set}[アクティブワープ集合／待機ワープ集合]
  pair.  Only the few active warps are searched; a warp that meets a
  long-latency operation moves to the pending set, and a pending warp
  whose operation has completed moves back to the active set when there
  is room (\cref{fig:l07-two-level}).
\end{definition}

Two papers are cited for this idea, with different
motives.\caution{Lecture vs paper: Gebhart et al.\ \cite{gebhart2011energy}
introduce the active/pending split for \emph{energy} (with a register
file cache, 36\,\% less register file energy).  Narasiman et al.\
\cite{narasiman2011twolevel} split the warps into fetch groups for
\emph{performance}: 19.1\,\% together with large warps.}  Narasiman
et al.\ divide the warps into
\term{fetch group}[フェッチグループ]s, for example 32 warps into 4
groups of 8; one group has priority until all its warps stall, with
round-robin inside and across groups.  The groups thus reach their
long-latency operations at different times, the same staggering as in
\cref{ex:l07-rr-gto}, and some group always has work while another
waits for memory.

\subsection{Scheduling and cache behaviour}

The scheduler also decides which data must be in the L1 cache at the
same time.

\begin{definition}[Working set, thrashing and compulsory misses]
  \label{def:l07-working-set}
  The \term{working set}[ワーキングセット] of a program over an
  interval is the set of memory blocks it uses in that interval.
  The cache suffers \term{cache thrashing}[キャッシュスラッシング] when
  the working set exceeds it, so that blocks are evicted before they are
  reused and most accesses miss.  A \term{compulsory miss}[初期参照ミス]
  is the miss of the first access to a block, which no cache can
  avoid.
\end{definition}

A hit on a block that the same warp brought into the cache exploits
\term{intra-warp locality}[ワープ内局所性]; a hit on a block another
warp brought in exploits \term{inter-warp locality}[ワープ間局所性].
\margin{Rogers et al.\ classify the hits of their benchmarks this way:
in the highly cache-sensitive ones most of the reuse is intra-warp
\cite{rogers2012ccws}.}  When each
warp reuses its own data, as in the next example, round-robin is the
worst order.

\begin{example}
  \label{ex:l07-thrashing}
  Three warps each read their own block four times: warp $k$ reads
  block B$k$.  The L1 holds two blocks with LRU replacement (Lesson~6),
  and each access hits or misses before the next one issues.
  \begin{center}
    \small
    \begin{tabular}{@{}ll@{}}
      \toprule
      order & accesses (bold = miss) \\
      \midrule
      RR  & \textbf{B1 B2 B3 B1 B2 B3 B1 B2 B3 B1 B2 B3} \\
      GTO & \textbf{B1} B1 B1 B1 \textbf{B2} B2 B2 B2 \textbf{B3} B3 B3 B3 \\
      \bottomrule
    \end{tabular}
  \end{center}
  Under round-robin each block is evicted by the two others before
  its warp returns to it: 12 misses in 12 accesses.  Under GTO each
  warp finishes its block before the next starts: only the 3
  compulsory misses.\caution{The GTO row assumes that a warp keeps
issuing after its first miss.  When other warps are ready at that miss,
GTO switches, and its working set holds more than one warp's data;
hence $\mathit{SW} \ge 1$ for GTO in the formula below.}
\end{example}

In general, let each warp have a working set of size $\mathit{SS}_W$
and let $N$ warps be resident on the SM.  Under round-robin all $N$
warps are in flight at once and the L1 must hold
$\mathit{SS}_W \times N$.  A scheduler that limits the number of
schedulable warps to $\mathit{SW}$ needs only
$\mathit{SS}_W \times \mathit{SW}$; for GTO, $\mathit{SW}$ is not
necessarily 1, because it switches warps at every long-latency event.
In \cref{ex:l07-thrashing}, $\mathit{SS}_W = 1$ block and $N = 3$, which
exceeds the two-block cache; with $\mathit{SW} = 2$ the working set
fits.\tip{L1 needed $\approx \mathit{SS}_W$ times the number of warps
that may issue: $N$ for round-robin, $\mathit{SW}$ under a limit.  To
fit a cache of size $C$, limit to $\mathit{SW} \le C/\mathit{SS}_W$.}

\begin{definition}[Cache-conscious wavefront scheduling]
  \label{def:l07-ccws}
  The scheduler of Rogers et al., called
  \term{cache-conscious wavefront scheduling (CCWS)}[キャッシュ考慮型ウェーブフロントスケジューリング]
  \cite{rogers2012ccws}, limits the number of warps that may issue so
  that their working set fits in the L1: on a stall it chooses among
  a small set of warps (W1--W3, say) and starts another warp (W4)
  only when those have completed.\index{CCWS|see{cache-conscious wavefront scheduling (CCWS)}}
\end{definition}

\caution{Lecture vs paper: the fixed limit described here is the
paper's baseline \enquote{static wavefront limiting}.  CCWS itself is
adaptive: a detector notices when evictions destroy a warp's
locality, and the scheduler then gives that warp more exclusive access
to the L1D by throttling the others.}%
The paper's mechanism is the
\term{lost intra-wavefront locality detector (LLD)}[ウェーブフロント内局所性損失検出器],
which keeps, for every warp, the tags of the lines that warp brought
into the L1 and lost to eviction (a victim tag array without data).
An L1 miss that hits in the warp's own victim tags is lost
intra-warp locality; it raises the warp's lost-locality score, and
warps pushed past a cutoff in the score stack (in the paper, above
it) by the high scores of others may not issue loads to the L1D until
the scores decay.  The
reported results are 25\,\% fewer L1 data cache misses on average and
a harmonic-mean speed-up of 24\,\% over earlier scheduling
policies.\margin{The CCWS hardware adds 0.17\,\% to the total chip
area; against a simple greedy scheduler the harmonic-mean gain on
the highly cache-sensitive benchmarks is 63\,\% \cite{rogers2012ccws}.}\index{LLD|see{lost intra-wavefront locality detector (LLD)}}

Limiting the warps raises two problems.
\begin{itemize}
  \item \emph{Synchronisation.}  A barrier (Lesson~3) releases the
    warps of a block only when all of them have arrived.  If the
    scheduler holds back a warp of a block whose other warps wait at
    \texttt{\_\_syncthreads()}, no warp of that block can progress: a
    deadlock.  The limit must therefore never keep a warp out that
    others are waiting for.
  \item \emph{The size of the limit.}  More warps give more
    thread-level parallelism to hide latency; fewer warps give a
    smaller working set.  The best limit depends on the cache and the
    program, which is why CCWS adapts it at run time.
\end{itemize}
\Needspace{8\baselineskip}
In a kernel of the following shape every warp of the block must reach
the barrier before any warp starts phase~2:

\begin{lstlisting}[language=C++, numbers=none]
__global__ void twoPhase(float *a) {
  // phase 1: every warp of the block works
  __syncthreads();   // waits for ALL warps
  // phase 2
}
\end{lstlisting}

\subsection{Other scheduling policies}

The scheduler can also serve the memory system beyond the L1.\margin{GPUMech (Lesson~8) models the performance of both RR and GTO
analytically, with average errors of 13.2\,\% and 14.0\,\%
\cite{huang2014gpumech}.}  Like the memory channels of Lesson~4, the
banks inside each DRAM can serve requests in parallel; the number of banks busy at the same
time is its \term{bank-level parallelism}[バンクレベル並列性].
\begin{itemize}
  \item \emph{Prefetch-aware scheduling} \cite{jog2013orchestrated}
    builds on the two-level scheduler but forms its groups from
    \emph{non-consecutive} warps.  Consecutive warps access
    neighbouring data, so a simple prefetcher triggered by one warp can
    fetch the data of the next; scheduled apart, the next warp runs
    late enough for the prefetch to arrive in time, and the requests in
    flight at once spread over more DRAM banks.\margin{Jog et al.\
report 25\,\% and 7\,\% average gains over round-robin and two-level
scheduling, each with prefetching, on a simulated 30-core GPU
\cite{jog2013orchestrated}.}
  \item \term{CTA-aware scheduling}[CTA考慮スケジューリング] (OWL
    \cite{jog2013owl}) decides which CTAs (thread blocks, Lesson~3) are
    scheduled together, grouping the warps of a few CTAs so that their
    data stays in the cache, and combines this with locality- and
    bank-level-parallelism-aware warp selection.\margin{OWL: a 33\,\%
average gain over round-robin warp scheduling on a simulated 28-core
GPU; in its motivation, 19 of 38 applications spend on average 62\,\%
of their cycles with no warp executing \cite{jog2013owl}.}
\end{itemize}

\begin{keypoint}[Lesson summary]
  GPU virtual memory translates warps, not threads: coalescing decides
  the number of translations, uncoalesced loads can need one per
  thread, and shared high-bandwidth page table walkers with page walk
  caches, multi-level TLBs, segment mapping and large pages keep the
  cost down; large pages trade TLB reach for internal fragmentation.
  A shared virtual address space replaces copy-then-execute by
  on-demand paging, whose far faults cost tens of microseconds over
  PCIe; tree-based neighbourhood prefetching amortises them, and
  eviction under oversubscription must respect the prefetcher's units.
  The warp scheduler issues in order within a warp but chooses among
  ready warps: round-robin advances all warps together and bunches
  their misses, GTO staggers them; two-level schedulers search a small
  active set for energy and performance; and limiting the schedulable
  warps (CCWS) shrinks the L1 working set, provided barriers are
  respected.
\end{keypoint}

\end{document}
